\SourcePage{220}
\SGAsection{8.4}{例：局部拓扑斯}
\textbf{8.4.1.} 设 $X=(X_i)_{i\in I}$ 是\Term{category}{范畴} $I$ 上的一个\Term{fibered-topos}{纤维拓扑斯}。对于 $I$ 的每个\Term{pro-object}{pro-对象} $\alpha=(i_\alpha)$，由公式
\[
 X_\alpha=\mathop{\varprojlim\nolimits^{\mathrm{top}}}_{I_\alpha}X_{i_\alpha}
 \tag{8.4.1.1}
\]
可得到\Term{category}{范畴} $\operatorname{Pro}(I)$（第 I 讲 8.10）上的一个\Term{fibered-topos}{纤维拓扑斯}（为说明这一点，可把第 I 讲习题 8.2.8 推广到\Term{pseudo-functor}{伪函子}）。

\textbf{8.4.2.} 设 $E$ 是一个\Term{topos}{拓扑斯}，$p:\operatorname{Fl}(E)\to E$ 是 7.3.1 中考虑的\Term{fibered-topos}{纤维拓扑斯}。如 8.4.1 那样延拓，可得到\Term{fibered-topos}{纤维拓扑斯} $\operatorname{Fl}(E)\to\operatorname{Pro}(E)$。\Term{category}{范畴} $\operatorname{Point}(E)$ 通过一个\Term{fully-faithful}{全忠实函子}映入 $\operatorname{Pro}(E)$（第 IV 讲 6.8.5），因而经\Term{base-change}{基变换}（7.1.9）得到\Term{fibered-topos}{纤维拓扑斯} $\operatorname{Loc}(E)\to\operatorname{Point}(E)$。在 $E$ 的点 $p$ 处的\Term{fiber}{纤维} $\operatorname{Loc}_p(E)$ 称为 $E$ 在点 $p$ 处的\Term{localized-topos}{局部化拓扑斯}。由上述构造，这个\Term{topos}{拓扑斯}\Term{covariant}{协变}地依赖于点 $p$。按定义，
\[
 \operatorname{Loc}_p(E)=\mathop{\varprojlim\nolimits^{\mathrm{top}}}_{X\in\operatorname{Vois}(p)}E_{/X},
 \tag{8.4.2.1}
\]
其中 $\operatorname{Vois}(p)$ 为 $p$ 的\Term{neighborhood}{邻域}\Term{category}{范畴}（第 IV 讲 6.8）。

\textbf{8.4.3.} 设 $m:p'\to p$ 是 $\operatorname{Point}(E)$ 中的一个\Term{morphism}{态射}（第 IV 讲 6）。对每个 $X\in\operatorname{Vois}(p)$，由此得到 $E_{/X}$ 的一个点 $p'_X$，再由 $\varprojlim^{\mathrm{top}}$ 的\Term{universal-property}{泛性质} 8.1，得到 $\operatorname{Loc}_p(E)$ 的一个点 $\theta_p(m)$。这样便定义了一个\Term{functor}{函子}
\[
 \theta_p:\operatorname{Point}(E)_{/p}\longrightarrow
 \operatorname{Point}(\operatorname{Loc}_p(E)),\tag{8.4.3.1}
\]
立即可见它是\Term{equivalence-of-categories}{范畴等价}。于是 $\operatorname{Point}(\operatorname{Loc}_p(E))$ \Term{canonical}{典范}\Term{equivalence}{等价}于 $p$ 的\Term{generization}{一般化}所成的\Term{category}{范畴}（第 IV 讲 7.1.8）。此外还有\Term{topos}{拓扑斯}的\Term{canonical}{典范}\Term{equivalence}{等价}
\[
 \mathrm{can}^{\bullet}:\operatorname{Loc}_{\theta(m)}(\operatorname{Loc}_p(E))
 \longrightarrow\operatorname{Loc}_{p'}(E),
\]
它包含在下列在\Term{isomorphism}{同构}意义下交换的图中：
\[
\begin{tikzcd}[column sep=large]
\operatorname{Loc}_{\theta(m)}(\operatorname{Loc}_p(E))
 \arrow[r,"\mathrm{can}^{\bullet}"]\arrow[dr,"\varphi"']
 &\operatorname{Loc}_{p'}(E)\arrow[d,"m^{\bullet}"]\\
 &\operatorname{Loc}_p(E).
\end{tikzcd}
\]
这里 $\varphi$ 是\Term{canonical}{典范}\Term{morphism}{态射}（8.4.2.1）。

\SourcePage{221}
\textbf{8.4.4.} \Term{localized-topos}{局部化拓扑斯}似乎只在来自\Term{algebra}{代数}几何的\Term{topos}{拓扑斯}情形，或至少在具有适当\Term{finite}{有限}性性质的\Term{topos}{拓扑斯}情形才有意义。例如，当 $E$ 是一个\Term{hausdorff-space}{分离拓扑空间}上的\Term{sheaf}{层}\Term{topos}{拓扑斯}时，立即可见（例如借助 8.2.9），\Term{localized-topos}{局部化拓扑斯}都是\Term{point-topos}{点拓扑斯}。另一方面，此时\Term{category}{范畴} $\operatorname{Point}(E)$ 是\Term{discrete}{离散}的，8.4.2 所述情形是平凡的。相反，设 $E=\operatorname{Top}(X)$ 是\Term{scheme}{概形} $X$ 上 Zariski \Term{topology}{拓扑}的\Term{sheaf}{层}\Term{topos}{拓扑斯}，$x$ 是 $X$ 的一点，$\mathscr O_x$ 是 $X$ 在 $x$ 处的\Term{local-ring}{局部环}，$Y=\operatorname{Spec}(\mathscr O_x)$。可以看出
\[
 \operatorname{Loc}_x(E)=\operatorname{Top}(Y).
\]
关于来自\Term{etale-topology}{平展拓扑}的其他例子，参见本研讨班第 VIII 讲。

\textbf{8.4.5.} 当 $E$ \Term{locally-coherent}{局部凝聚}（2.3）时，\Term{localized-topos}{局部化拓扑斯} $\operatorname{Loc}_p(E)$ 是\Term{coherent}{凝聚的}（在（8.4.2.1）中可只考虑既\Term{algebraic}{代数}又\Term{coherent}{凝聚}的 $X\in\operatorname{Vois}(p)$，然后应用 8.3.13）。对每个点\Term{morphism}{态射} $m:p'\to p$，\Term{topos}{拓扑斯}\Term{morphism}{态射} $m^{\bullet}:\operatorname{Loc}_{p'}(E)\to\operatorname{Loc}_p(E)$ 是\Term{coherent-morphism}{凝聚态射}（8.3.14）。

\textbf{8.4.6.} 如果\Term{topos}{拓扑斯} $X$ 的\Term{functor}{函子} $\Gamma(X,-)=$“$X$ 上的\Term{section}{截面}”（第 IV 讲 4.3）是一个\Term{fiber-functor}{纤维函子}（第 IV 讲 6），就称 $X$ 为\Term{local-topos}{局部拓扑斯}。与该\Term{fiber-functor}{纤维函子}对应的点称为\Term{local-topos}{局部拓扑斯}的\Term{center}{中心}。当 $E$ \Term{locally-coherent}{局部凝聚}时，\Term{localized-topos}{局部化拓扑斯}是\Term{local-topos}{局部拓扑斯}（1.2.3、8.5.2 和 8.5.7）。$\operatorname{Loc}_p(E)$ 的\Term{center}{中心}\Term{canonical-isomorphism}{典范同构}于 $\theta_p(\operatorname{id}_p)$（8.4.3.1）。$\operatorname{Loc}_p(E)$ 的\Term{center}{中心}在 $E$ 中的像\Term{isomorphism}{同构}于 $p$。配备\Term{canonical}{典范}\Term{morphism}{态射} $\operatorname{Loc}_p(E)\to E$ 的\Term{localized-topos}{局部化拓扑斯} $\operatorname{Loc}_p(E)$，是把\Term{local-topos}{局部拓扑斯}映入 $E$、并把\Term{center}{中心}映“到”$p$ 上这一\Term{universal-problem}{泛问题}（二维\Term{universal-problem}{泛问题}！）的解。

\textbf{8.4.7.} 关于\Term{local-topos}{局部拓扑斯}，还有若干问题，撰写者尚未研究。例如，若 $X$ 为\Term{local-topos}{局部拓扑斯}，则 $X$ 的\Term{final-object}{终对象} $e$ 有一个极大的开\Term{subobject}{子对象} $U\ne e$。$U$ 的补 $Y$ 是一个没有非平凡开\Term{subobject}{子对象}的\Term{local-topos}{局部拓扑斯}。这样的\Term{topos}{拓扑斯}是否为\Term{point-topos}{点拓扑斯}？设 $U$ 为\Term{topos}{拓扑斯}，$X=(U,\Phi:U\to(\mathrm{Ens}))$ 是通过\Term{gluing}{粘合}得到的\Term{topos}{拓扑斯}。\Term{gluing-functor}{粘合函子} $\Phi$ 满足什么条件时 $X$ 才是\Term{local-topos}{局部拓扑斯}？

\SGAsection{8.5}{拓扑斯的滤过投射极限上层的结构}
\textbf{8.5.1.} 本节中，$F=(F_i)_{i\in I}$ 是小\Term{cofiltered-category}{余滤过范畴} $I$ 上的一个 $\mathcal U$-\Term{fibered-topos}{纤维拓扑斯}。选择 $F$ 的一个\Term{bicleavage}{双分裂}（7.1.4），即对每个 $f\in\operatorname{Fl}(I):i\to j$，选择\Term{topos}{拓扑斯}\Term{morphism}{态射} $f_\bullet:F_i\to F_j$，使 $f^*:F_j\to F_i$ 为\Term{fibered-structure}{纤维结构}的\Term{inverse-image-functor}{逆像函子}。于是有满足\Term{cocycle-condition}{余圈条件}的\Term{canonical-isomorphism}{典范同构} $c_{f,g}$（7.1.3）。以 $\underleftarrow F$ 表示\Term{fibered-topos}{纤维拓扑斯} $F$ 的\Term{projective-limit}{投射极限}（8.2.3），对每个 $i\in\operatorname{ob}I$，以
\[
 \mu_i:\underleftarrow F\longrightarrow F_i\tag{8.5.1.1}
\]
\SourcePage{222}
表示\Term{canonical}{典范}\Term{morphism}{态射}（8.1.3）。记 $\operatorname{Top}(F)$ 为 $F$ 的\Term{total-topos}{全拓扑斯}（7.4.3，3）。由 8.2.9，\Term{canonical-functor}{典范函子} $F\to\underrightarrow F$ 定义一个\Term{topos}{拓扑斯}\Term{morphism}{态射}
\[
 Q:\underleftarrow F\longrightarrow\operatorname{Top}(F).\tag{8.5.1.2}
\]
\Term{topos}{拓扑斯} $\underleftarrow F$ 可等同于 $\mathscr{H}\!om_{\mathrm{Cart}/I^\circ}(I^\circ,F')$（8.2.9），而 $\operatorname{Top}(F)$ 可等同于 $\mathscr{H}\!om_{I^\circ}(I^\circ,F')$（7.4.7）。作这些等同后，\Term{morphism}{态射} $Q$ 就是 $\mathscr{H}\!om_{\mathrm{Cart}/I^\circ}(I^\circ,F')$ 到 $\mathscr{H}\!om_{I^\circ}(I^\circ,F')$ 的\Term{embedding}{嵌入}\Term{morphism}{态射}（8.2.9）；对 $\underleftarrow F$ 的每个\Term{object}{对象} $M$，有
\[
 Q_*(M)=(i\longmapsto\mu_{i*}(M)).\tag{8.5.1.3}
\]
由 7.4.3.4，对每个 $i\in\operatorname{ob}I$，有\Term{topos}{拓扑斯}\Term{morphism}{态射}
\[
 \alpha_i:F_i\longrightarrow\operatorname{Top}(F).\tag{8.5.1.4}
\]
图
\[
\begin{tikzcd}
\underleftarrow F\arrow[r,"\mu_i"]\arrow[dr,"Q"']&F_i\arrow[d,"\alpha_i"]\\
&\operatorname{Top}(F)
\end{tikzcd}\tag{8.5.1.5}
\]
一般并不交换（即使在\Term{isomorphism}{同构}意义下也是如此）。然而，由有关\Term{morphism}{态射}的定义，有\Term{canonical-isomorphism}{典范同构}
\[
 \alpha_i^*Q_*\simeq\mu_{i*}.\tag{8.5.1.6}
\]
当 $i$ 是 $I$ 的\Term{final-object}{终对象}时，$\alpha_i$ 是一个有\Term{retraction}{收缩} $\beta_i:\operatorname{Top}(F)\to F_i$ 的\Term{embedding}{嵌入}（7.4.12），而图
\[
\begin{tikzcd}
\underleftarrow F\arrow[r,"\mu_i"]\arrow[dr,"Q"']&F_i\\
&\operatorname{Top}(X)\arrow[u,"\beta_i"']
\end{tikzcd}\tag{8.5.1.7}
\]
在\Term{canonical-isomorphism}{典范同构}意义下交换。

\begin{SGAresult}{命题}{8.5.2}
设 $j\mapsto M_j$ 为 $\operatorname{Top}(F)$ 的一个\Term{object}{对象}。存在\Term{functorial}{函子性}\Term{isomorphism}{同构}
\[
 Q^*(j\mapsto M_j)\simeq\varinjlim_{I^\circ}\mu_j^*M_j.\tag{8.5.2.1}
\]
\end{SGAresult}
\begin{proof}
$\operatorname{Top}(F)$ 的一个\Term{object}{对象}（7.1.3）由一个映射 $j\mapsto M_j$、$M_j\in\operatorname{ob}F_j$，以及对每个\Term{morphism}{态射} $f:i\to j$ 给定的\Term{morphism}{态射}
\[
 \beta_f:M_j\longrightarrow f_*M_i\tag{8.5.2.2}
\]
\SourcePage{223}
组成；\Term{equivalence}{等价}地，借助\Term{adjunction}{伴随}，可给定\Term{morphism}{态射}
\[
 \beta'_f:f^*M_j\longrightarrow M_i,\tag{8.5.2.3}
\]
这些\Term{morphism}{态射}须满足：对于每对可复合\Term{morphism}{态射} $i\xrightarrow f j\xrightarrow g k$，下图交换：
\[
\begin{tikzcd}
f^*g^*M_k\arrow[r,"f^*\beta'_g"]\arrow[d,"{c^*_{g,f}}"']&f^*M_j\arrow[d,"\beta'_f"]\\
(gf)^*M_k\arrow[r,"\beta'_{gf}"']&M_i.
\end{tikzcd}\tag{8.5.2.4}
\]
再回顾（8.1.3.1），对于每个\Term{morphism}{态射} $f:i\to j$，有\Term{isomorphism}{同构} $b_f:f_\bullet\mu_i\xrightarrow\sim\mu_j$；对于每对可复合\Term{morphism}{态射} $i\xrightarrow f j\xrightarrow g k$，有\Term{commutative-diagram}{交换图}
\[
\begin{tikzcd}
g_\bullet f_\bullet\mu_i\arrow[r,"{c_{g,f}}"]\arrow[d,"g_\bullet(b_f)"']&(gf)_\bullet\mu_i\arrow[d,"b_{gf}"]\\
g_\bullet\mu_j\arrow[r,"b_g"']&\mu_k.
\end{tikzcd}\tag{8.5.2.5}
\]
现设 $f:i\to j$ 为 $I$ 中的\Term{morphism}{态射}。以
\[
 t_f:\mu_j^*(M_j)\longrightarrow\mu_i^*(M_i)\tag{8.5.2.6}
\]
表示复合\Term{morphism}{态射}
\[
 \mu_j^*(M_j)\xrightarrow{b_f^*}\mu_j^*(f^*(M_j))
 \xrightarrow{\mu_i^*(\beta'_f)}\mu_i^*(M_i).
\]
由图（8.5.2.4）和（8.5.2.5）的交换性可得，对每对可复合\Term{morphism}{态射} $i\to j\to k$，有
\[
 t_ft_g=t_{gf}.\tag{8.5.2.7}
\]
其核实留给读者。于是定义了\Term{functor}{函子} $I^\circ\to\underleftarrow F$，可考虑其\Term{inductive-limit}{归纳极限} $\varinjlim_{I^\circ}\mu_i^*(M_i)$。设 $N$ 为 $\underleftarrow F$ 的\Term{object}{对象}。有
\[
 \operatorname{Hom}(\varinjlim_{I^\circ}\mu_i^*(M_i),N)
 \simeq\varprojlim_{I^\circ}\operatorname{Hom}(\mu_i^*(M_i),N),\tag{8.5.2.8}
\]
从而由\Term{adjunction}{伴随}，有\Term{isomorphism}{同构}
\[
 \operatorname{Hom}(\varinjlim_{I^\circ}\mu_i^*(M_i),N)
 \simeq\varprojlim_{I^\circ}\operatorname{Hom}_{F_i}(M_i,\mu_{i*}(N)).\tag{8.5.2.9}
\]
我们要解释（8.5.2.9）的右端。为此，对每个\Term{morphism}{态射} $f:i\to j$，以
\[
 d_f:\operatorname{Hom}_{F_i}(M_i,\mu_{i*}(N))
 \longrightarrow\operatorname{Hom}_{F_j}(M_j,\mu_{j*}(N))
\]
\SourcePage{224}
表示（8.5.2.9）中\Term{projective-system}{投射系统}的\Term{transition-morphism}{转移态射}。由\Term{morphism}{态射} $t_f$ 的定义（8.5.2.6），映射 $d_f$ 把\Term{morphism}{态射} $u_i:M_i\to\mu_{i*}(N)$ 送到\Term{morphism}{态射} $d_f(u_i):M_j\to\mu_{j*}(N)$，后者由以下\Term{morphism}{态射}复合而得：
\[
 M_j\xrightarrow{\beta_f}f_*(M_i)\xrightarrow{f_*(u_i)}f_*\mu_{i*}(N)
 \xrightarrow{b_f}\mu_{j*}(N).
\]
因此，（8.5.2.9）右端的元素\Term{solvable}{可解}释为一族\Term{morphism}{态射} $u_i:M_i\to\mu_{i*}(N)$，$i\in\operatorname{ob}I$，使对每个 $f:i\to j$，图
\[
\begin{tikzcd}
M_j\arrow[r]\arrow[d,"u_j"']&f_*(M_i)\arrow[d,"f_*(u_i)"]\\
\mu_{j*}(N)\arrow[r,"\sim"']&f_*\mu_{i*}(N)
\end{tikzcd}
\]
交换，也就是从\Term{section}{截面} $(i\mapsto M_i)\in\mathscr{H}\!om_{I^\circ}(I^\circ,F')$ 到\Term{section}{截面} $Q_*(N)$ 的一个\Term{morphism}{态射}。因此有\Term{functorial}{函子性}\Term{isomorphism}{同构}
\[
 \operatorname{Hom}(\varinjlim_{I^\circ}\mu_{i*}(M_i),N)
 \simeq\operatorname{Hom}((i\mapsto M_i),Q_*(N));
\]
再由\Term{adjunction}{伴随}即得命题。
\end{proof}

\begin{SGAresult}{命题}{8.5.3}
若对每个 $f:i\to j\in\operatorname{Fl}(I)$，\Term{functor}{函子} $f_*:F_i\to F_j$ 与小\Term{filtered-inductive-limit}{滤过归纳极限}可交换，则对每个\Term{section}{截面} $(i\mapsto M_i)\in\mathscr{H}\!om_{I^\circ}(I^\circ,F')$，有
\[
 Q_*Q^*(i\mapsto M_i)\simeq
 \left(i\longmapsto\varinjlim_{f:j\to i}f_*M_j\right).\tag{8.5.3.1}
\]
\end{SGAresult}
\begin{proof}
用 $\beta_f$（8.5.2.2）和 $c_{f,g}$（8.5.1）具体写出\Term{inductive-system}{归纳系统} $(f:j\to i)\mapsto f_*(M_j)$ 的\Term{transition-morphism}{转移态射}，以及\Term{section}{截面} $i\mapsto\varinjlim_{f:j\to i}f_*M_i$ 的\Term{transition-morphism}{转移态射}，均留给读者。由于\Term{functor}{函子} $f_*$ 与\Term{filtered-inductive-limit}{滤过归纳极限}可交换，\Term{section}{截面} $i\mapsto\varinjlim_{f:j\to i}f_*M_j$ 是\Term{cartesian-section}{笛卡尔截面}。又有从 $(i\mapsto M_i)$ 到\Term{section}{截面} $(i\mapsto\varinjlim_{f:j\to i}f_*M_j)$ 的自然\Term{morphism}{态射} $u$，显然从 $(i\mapsto M_i)$ 到一个\Term{cartesian-section}{笛卡尔截面}的每个\Term{morphism}{态射}都唯一地经 $u$ \Term{resolution}{分解}，由此得到\Term{isomorphism}{同构}（8.5.3.1）。
\end{proof}

\begin{SGAresult}{推论}{8.5.4}
在 8.5.3 的假设下，\Term{functor}{函子} $Q_*$ 与小\Term{filtered-inductive-limit}{滤过归纳极限}可交换。
\end{SGAresult}
\SourcePage{225}
\begin{proof}
设 $\alpha\mapsto N^\alpha$ 是 $\underleftarrow F$ 的一个小\Term{filtered-inductive-system}{滤过归纳系统}。置 $M^\alpha=Q_*(N^\alpha)$，故 $M^\alpha$ 是\Term{section}{截面} $i\mapsto M_i^\alpha$。有 $N^\alpha\simeq Q^*Q_*(N^\alpha)$；由于 $Q^*$ 与\Term{inductive-limit}{归纳极限}可交换，得到 $\varinjlim_\alpha N^\alpha\simeq Q^*(\varinjlim_\alpha M^\alpha)$。$\operatorname{Top}(F)=\mathscr{H}\!om_{I^\circ}(I^\circ,F')$ 中的\Term{inductive-limit}{归纳极限}逐\Term{fiber}{纤维}计算。因此 $\varinjlim_\alpha M\simeq(i\mapsto\varinjlim_\alpha M_i^\alpha)$。由（8.5.3.1），
\[
 Q_*(\varinjlim_\alpha N^\alpha)\simeq Q_*Q^*(\varinjlim_\alpha M^\alpha)
 \simeq\left(i\longmapsto\varinjlim_{f:j\to i}f_*(\varinjlim_\alpha M_j^\alpha)\right).
\]
因 $f_*$ 与\Term{filtered-inductive-limit}{滤过归纳极限}可交换，得到
\[
\begin{split}
 Q_*(\varinjlim_\alpha N^\alpha)
 &\simeq\left(i\longmapsto\varinjlim_{f:j\to i}\varinjlim_\alpha f_*(M_j^\alpha)\right)\\
 &\simeq\left(i\longmapsto\varinjlim_\alpha\varinjlim_{f:j\to i}f_*(M_j^\alpha)\right)
 \simeq\varinjlim_\alpha Q_*(N^\alpha).
\end{split}
\]
\end{proof}

\begin{SGAresult}{推论}{8.5.5}
若\Term{functor}{函子} $f_*:F_i\to F_j$ 与小\Term{filtered-inductive-limit}{滤过归纳极限}可交换，则对 $I$ 的每个\Term{object}{对象} $i$ 和每个 $M_i\in\operatorname{ob}F_i$，有
\[
 \mu_{i*}\mu_i^*(M_i)\simeq\varinjlim_{f:j\to i}f_*f^*(M_i).\tag{8.5.5.1}
\]
\end{SGAresult}
\begin{proof}
作\Term{base-change}{基变换} $I_{/i}\to I$，可设 $i$ 是 $I$ 的\Term{final-object}{终对象}（8.2.1）。此时\Term{morphism}{态射} $\mu_i:\underleftarrow F\to F_i$ 是以下\Term{morphism}{态射}的复合（8.2.9 和 7.4.12）：
\[
 Q:\underleftarrow F\longrightarrow F^\sim,\qquad
 \beta_i:(\alpha_{i!},\alpha_i^*):F^\sim\longrightarrow F_i.
\]
故 $\mu_i^*(M_i)\simeq Q^*\beta_i^*(M_i)$，这里 $\beta_i^*(M_i)$ 是\Term{section}{截面} $(f:i\to j)\mapsto f^*(M_i)$，其\Term{transition-morphism}{转移态射}由 $c_{f,g}$ \Term{derivation}{导出}。断言随即由 8.5.3 得出。
\end{proof}

\begin{SGAresult}{推论}{8.5.6}
在 8.5.5 的假设下，\Term{functor}{函子} $\mu_{i*}:\underleftarrow F\to F_i$ 与\Term{filtered-inductive-limit}{滤过归纳极限}可交换。
\end{SGAresult}
\begin{proof}
$\mu_{i*}$ 是\Term{functor}{函子} $Q_*$ 与“\Term{restriction}{限制}到 $i$ 处\Term{fiber}{纤维}”的\Term{functor}{函子}的复合；前者与\Term{filtered-inductive-limit}{滤过归纳极限}可交换，后者与\Term{inductive-limit}{归纳极限}可交换。
\end{proof}

\begin{SGAresult}{推论}{8.5.7}
在 8.5.3 的假设下，设 $i$ 是 $I$ 的\Term{object}{对象}，$X$ 是 $F_i$ 的\Term{object}{对象}，且 $F_i$ 上的\Term{functor}{函子} $\operatorname{Hom}(X,-)$ 与\Term{filtered-inductive-limit}{滤过归纳极限}可交换。则\Term{functor}{函子} $\operatorname{Hom}(\mu_i^*(X),-)$ 与\Term{filtered-inductive-limit}{滤过归纳极限}可交换；对 $F_i$ 的每个\Term{object}{对象} $M_i$，有
\[
 \operatorname{Hom}(\mu_i^*(X),\mu_i^*(M_i))\simeq
 \varinjlim_{f:j\to i}\operatorname{Hom}(f^*(X),f^*(M_i)).\tag{8.5.7.1}
\]
\SourcePage{226}
对 $\operatorname{Top}(F)$ 的每个\Term{object}{对象} $(j\mapsto M_j)$，有
\[
 \operatorname{Hom}(\mu_i^*(X),Q^*(j\mapsto M_j))\simeq
 \varinjlim_{f:j\to i}\operatorname{Hom}(f^*(X),M_j).\tag{8.5.7.2}
\]
\end{SGAresult}
\begin{proof}
\Term{functor}{函子} $\operatorname{Hom}(\mu_i^*(X),-)$ \Term{isomorphism}{同构}于 $\operatorname{Hom}(X,\mu_{i*}(-))$，后者与\Term{filtered-inductive-limit}{滤过归纳极限}可交换（8.5.4）。由（8.5.5.1），有
\[
 \operatorname{Hom}(\mu_i^*(X),\mu_i^*(M_i))
 \simeq\operatorname{Hom}(X,\mu_{i*}\mu_i^*(M_i))
 \simeq\operatorname{Hom}(X,\varinjlim_{f:j\to i}f_*f^*(M_i)),
\]
由此得到（8.5.7.1）。同样，由（8.5.3.1），有
\[
\begin{split}
 \operatorname{Hom}(\mu_i^*(X),Q^*(j\mapsto M_j))
 &\simeq\operatorname{Hom}(X,\mu_{i*}Q^*(j\mapsto M_j))\\
 &\simeq\operatorname{Hom}(X,\varinjlim_{f:j\to i}f_*(M_j)),
\end{split}
\]
由此得到（8.5.7.2）。
\end{proof}

\begin{SGAresult}{命题}{8.5.8}
设 $G\to I$ 为一个 $\mathcal U$-\Term{fibered-topos}{纤维拓扑斯}，$m:F\to G$ 为\Term{fibered-topos}{纤维拓扑斯}的\Term{cartesian-morphism}{笛卡尔态射}（7.1.15）。对 $G$ 使用 8.1.1 引入的记号。设 $\underleftarrow m:\underleftarrow F\to\underleftarrow G$ 为对 $m$ 取\Term{projective-limit}{投射极限}而得的\Term{morphism}{态射}（8.1.4）。下面由\Term{topos}{拓扑斯}与\Term{topos}{拓扑斯}\Term{morphism}{态射}组成的图在\Term{isomorphism}{同构}意义下交换：
\[
\begin{tikzcd}
\underleftarrow F\arrow[r,"Q"]\arrow[d,"\underleftarrow m"']&\operatorname{Top}(F)\arrow[d,"m^\sim"]\\
\underleftarrow G\arrow[r,"Q"']&\operatorname{Top}(G).
\end{tikzcd}\tag{8.5.8.1}
\]
对 $\underleftarrow F$ 的每个\Term{object}{对象} $N$，有
\[
 (\underleftarrow m)_*(N)\simeq
 \varinjlim_{I^\circ}\mu_j^*m_{j*}\mu_{j*}(N).\tag{8.5.8.2}
\]
\end{SGAresult}
\begin{proof}
图（8.5.8.1）的交换性直接来自定义（8.2.8 和 8.5.1）。有 $(\underleftarrow m)_*(N)\simeq Q^*Q_*(\underleftarrow m)_*(N)$；由（8.5.8.1）的交换性，$(\underleftarrow m)_*(N)\simeq Q^*m_*^\sim Q_*(N)$。所以有\Term{isomorphism}{同构} $(\underleftarrow m)_*(N)\simeq Q^*m_*^\sim(j\mapsto\mu_{j*}(N))$。\Term{functor}{函子} $m_*^\sim$ 就是 $\mathscr{H}\!om_{I^\circ}(I^\circ,m_*)$（7.4.10），因而有\Term{canonical-isomorphism}{典范同构} $(\underleftarrow m)_*(N)\simeq Q^*(j\mapsto m_{j*}\mu_{j*}(N))$。公式（8.5.8.2）于是由 8.5.2 得出。
\end{proof}

\begin{SGAresult}{命题}{8.5.9}
在 8.5.8 的假设和记号下，再假定对 $I$ 的每个\Term{object}{对象} $i$，\Term{functor}{函子} $m_{i*}:F_i\to G_i$ 与\Term{filtered-inductive-limit}{滤过归纳极限}可交换；对 $I$ 的每个\Term{morphism}{态射} $f:i\to j$，\Term{functor}{函子} $f_*:F_i\to F_j$ 与\Term{filtered-inductive-limit}{滤过归纳极限}可交换。则\Term{functor}{函子} $(\underleftarrow m)_*:\underleftarrow F\to\underleftarrow G$ 与\Term{filtered-inductive-limit}{滤过归纳极限}
\SourcePage{227}
可交换；对 $\operatorname{Top}(F)$ 的每个\Term{object}{对象} $(i\mapsto M_i)$，有\Term{canonical-isomorphism}{典范同构}
\[
 (\underleftarrow m)_*Q^*(i\mapsto M_i)\simeq
 \varinjlim_{I^\circ}\mu_j^*m_{j*}(M_j).\tag{8.5.9.1}
\]
\end{SGAresult}
\begin{proof}
第一个断言由（8.5.8.2）和 8.5.6 得出。由（8.5.8.2）与（8.5.3.1），有\Term{functorial}{函子性}\Term{isomorphism}{同构}
\[
 (\underleftarrow m)_*Q^*(i\mapsto M_i)\simeq
 \varinjlim_i\mu_i^*m_{i*}(\varinjlim_{f:j\to i}f_*(M_j)).
\]
利用 $m_{i*}$ 与\Term{filtered-inductive-limit}{滤过归纳极限}的交换性，以及\Term{isomorphism}{同构} $m_{i*}f_*\simeq f_*m_{j*}$（7.1.6），得到
\[
 (\underleftarrow m)_*Q^*(i\mapsto M_i)\simeq
 \varinjlim_i\mu_i^*(\varinjlim_{f:j\to i}f_*m_{j*}(M_j)).
\]
因 $\mu_i^*$ 与\Term{inductive-limit}{归纳极限}可交换，得到
\[
 (\underleftarrow m)_*Q^*(i\mapsto M_i)\simeq
 \varinjlim_i\varinjlim_{f:j\to i}\mu_i^*f_*m_{j*}(M_j).
\]
最后这个\Term{object}{对象}\Term{solvable}{可解}释为\Term{category}{范畴} $\operatorname{Fl}(I)^\circ$ 上的\Term{inductive-limit}{归纳极限}，其中 $\operatorname{Fl}(I)$ 是 $I$ 的\Term{morphism}{态射}\Term{category}{范畴}。设 $\varphi:I\to\operatorname{Fl}(I)$ 为把 $I$ 的每个\Term{object}{对象} $i$ 送到 $i$ 的恒等\Term{morphism}{态射}的\Term{functor}{函子}。\Term{functor}{函子} $\varphi^\circ:I^\circ\to\operatorname{Fl}(I)^\circ$ 是\Term{cofinal}{共尾的}，因而（第 I 讲 8.1）有\Term{canonical-isomorphism}{典范同构}
\[
 (\underleftarrow m)_*Q^*(i\mapsto M_i)\simeq
 \varinjlim_{I^\circ}\mu_j^*m_{j*}(M_j).
\]
\end{proof}

\begin{SGAresult}{推论}{8.5.10}
在 8.5.9 的条件下，对 $I$ 的每个\Term{object}{对象} $i$ 与 $F_i$ 的每个\Term{object}{对象} $M_i$，有\Term{canonical-isomorphism}{典范同构}
\[
 (\underleftarrow m)_*\mu_i^*(M_i)\simeq
 \varinjlim_{f:j\to i}\mu_j^*m_{j*}f^*(M_i).\tag{8.5.10.1}
\]
\end{SGAresult}
\begin{proof}
证明与推论 8.5.5 的证明类似。
\end{proof}

\textbf{注记 8.5.11。} 回顾，\Term{coherent}{凝聚}\Term{topos}{拓扑斯}之间\Term{coherent-morphism}{凝聚态射}的\Term{direct-image-functor}{正像函子}与\Term{filtered-inductive-limit}{滤过归纳极限}可交换（5.1）。因此，当所考虑的\Term{fibered-topos}{纤维拓扑斯}是\Term{coherent}{凝聚的}，且所考虑的\Term{fibered-topos}{纤维拓扑斯}\Term{morphism}{态射}是\Term{coherent-morphism}{凝聚态射}时，8.5.3 至 8.5.7，以及 8.5.9、8.5.10 均适用。

\SGAsection{8.6}{赋环的纤维 $\mathcal U$-拓扑斯}
\SourcePage{228}
\textbf{8.6.1.} 若一个 $\mathcal U$-\Term{fibered-topos}{纤维拓扑斯}在其\Term{total-site}{全景}（7.4.1）上配有一个\Term{sheaf-of-rings}{环层}，便称它是\Term{ringed}{赋环的}。设 $p:F\to I$ 为一个 $\mathcal U$-\Term{fibered-topos}{纤维拓扑斯}。选择一个\Term{bicleavage}{双分裂}（7.1.4）。由 7.4.7，在 $F$ 上给定一个\Term{sheaf-of-rings}{环层}，\Term{equivalence}{等价}于对 $I$ 的每个\Term{object}{对象} $i$ 给定 $F_i$ 中的一个环 $A_i$，并对每个\Term{morphism}{态射} $f:i\to j$ 给定\Term{morphism}{态射} $\varphi_f:A_j\to f_*A_i$，且这族 $\varphi_f$ 须满足该处写明的相容条件。因此，\Term{topos}{拓扑斯}\Term{morphism}{态射} $f_\bullet:F_i\to F_j$ 是\Term{ringed-topos}{赋环拓扑斯}（分别赋以 $A_i$ 和 $A_j$）的\Term{morphism}{态射}（第 IV 讲 11），而 $F$ 上的\Term{ringed}{赋环结构}及\Term{morphism}{态射} $f_\bullet$ 的选择，给出从 $I$ 到\Term{ringed-topos}{赋环拓扑斯}的二维\Term{category}{范畴}的一个\Term{pseudo-functor}{伪函子}。反之，给定这样一个\Term{pseudo-functor}{伪函子}，便可重建出产生它的\Term{ringed}{赋环} $\mathcal U$-\Term{fibered-topos}{纤维拓扑斯}。

\textbf{8.6.2.} 如 8.1 那样，可以定义\Term{ringed}{赋环} $\mathcal U$-\Term{topos}{拓扑斯} $(F_i,A_i,i\in I)$ 的\Term{projective-limit}{投射极限}。我们假定这个直接的推广已经作出，并在这种情形自由使用 8.1 引入的语言。当此\Term{projective-limit}{投射极限}存在时，它是一个\Term{ringed}{赋环} $\mathcal U$-\Term{topos}{拓扑斯} $(H,B)$，且对 $I$ 的每个\Term{object}{对象} $i$，有\Term{ringed-topos}{赋环拓扑斯}\Term{morphism}{态射} $\mu_i:(H,B)\to(F_i,A_i)$。若（不考虑\Term{ringed}{赋环}的）\Term{fibered-topos}{纤维拓扑斯} $F$ 有\Term{projective-limit}{投射极限} $\varprojlim_I^{\mathrm{top}}F_i$，并且环的\Term{inductive-system}{归纳系统} $i\mapsto\mu_i^*(A_i)$ 在 $\varprojlim_I^{\mathrm{top}}F_i$ 的环\Term{category}{范畴}中有\Term{inductive-limit}{归纳极限} $\varinjlim_{I^\circ}\mu_i^*A_i$（$I$ 小时总是如此），则\Term{ringed-topos}{赋环拓扑斯}
\[
 \left(\mathop{\varprojlim\nolimits^{\mathrm{top}}}_I F_i,
 \varinjlim_{I^\circ}\mu_i^*A_i\right)
\]
显然是所考虑的\Term{ringed}{赋环}\Term{fibered-topos}{纤维拓扑斯}的\Term{projective-limit}{投射极限}。

\textbf{8.6.3.} 第 8.5 节为\Term{sheaf-of-sets}{集合层}建立的公式也适用于\Term{abelian-sheaf}{阿贝尔层}，并且也适用于\Term{sheaf-of-modules}{模层}，但其中的\Term{inverse-image-functor}{逆像函子}须解释为\Term{sheaf-of-modules}{模层}意义下的\Term{inverse-image-functor}{逆像函子}（第 IV 讲 1.1）。核实留给读者。

\textbf{8.6.4.} 设 $(p:F\to I,A)$ 是\Term{cofiltered-category}{余滤过范畴}上的\Term{ringed}{赋环} $\mathcal U$-\Term{fibered-topos}{纤维拓扑斯}。若对每个\Term{morphism}{态射} $f:i\to j$，\Term{morphism}{态射} $f_\bullet:(F_i,A_i)\to(F_j,A_j)$ 是\Term{ringed-topos}{赋环拓扑斯}的\Term{right-flat}{右平坦}（相应地，\Term{left-flat}{左平坦}）\Term{morphism}{态射}（第 V 讲 1.8），便称 $(p:F\to I,A)$ 为\Term{right-flat}{右平坦}（相应地，\Term{left-flat}{左平坦}）的\Term{ringed}{赋环} $\mathcal U$-\Term{fibered-topos}{纤维拓扑斯}。假定 $I$ \Term{essentially-small}{本质小}。则\Term{ringed-topos}{赋环拓扑斯}\Term{morphism}{态射}
\[
 Q:\left(\mathop{\varprojlim\nolimits^{\mathrm{top}}}_I F_i,
 \varinjlim_{I^\circ}\mu_i^*(A_i)\right)
 \longrightarrow(\mathscr{H}\!om_{I^\circ}(I^\circ,F'),A)
\]
（8.5.1）是\Term{right-flat}{右平坦的}（相应地，\Term{left-flat}{左平坦的}），这由定义与 8.5.2 可得，而且不需要对\Term{transition-morphism}{转移态射} $f_\bullet$ 作\Term{flat}{平坦性}假设。若再假定 $(p:F\to I,A)$ 是\Term{right-flat}{右平坦}（相应地，\Term{left-flat}{左平坦}）的\Term{ringed}{赋环} $\mathcal U$-\Term{fibered-topos}{纤维拓扑斯}，则\Term{morphism}{态射}
\[
 \mu_i:\left(\mathop{\varprojlim\nolimits^{\mathrm{top}}}_I F_i,
 \varinjlim_{I^\circ}\mu_i^*A_i\right)\longrightarrow(F_i,A_i)
\]
是\Term{right-flat}{右平坦的}（相应地，\Term{left-flat}{左平坦的}）。这来自\Term{flat-module}{平坦模}的\Term{inverse-image}{逆像}仍\Term{flat}{平坦}（第 V 讲 1.7.1），以及\Term{flat-module}{平坦模}的\Term{filtered-inductive-limit}{滤过归纳极限}仍为\Term{flat-module}{平坦模}。同样的论证还可证明：
\SourcePage{229}
若 $m:(F_i,A_i,i\in I)\to(G_i,A_i,i\in I)$ 是一个\Term{essentially-small}{本质小}的\Term{cofiltered-category}{余滤过范畴}上\Term{ringed}{赋环} $\mathcal U$-\Term{fibered-topos}{纤维拓扑斯}的\Term{cartesian-morphism}{笛卡尔态射}，且对每个\Term{object}{对象} $i\in\operatorname{ob}(I)$，\Term{morphism}{态射} $m_i$ 是\Term{right-flat}{右平坦的}（相应地，\Term{left-flat}{左平坦的}），则对 $m$ 取\Term{projective-limit}{投射极限}而得的\Term{morphism}{态射} $\underleftarrow m$ 是\Term{right-flat}{右平坦的}（相应地，\Term{left-flat}{左平坦的}）。

\SGAsection{8.7}{拓扑斯投射极限上层的上同调}
\textbf{8.7.1.} 本节中，$I$ 表示小\Term{cofiltered-category}{余滤过范畴}，$(p:F\to I,A)$ 与 $(q:G\to I,B)$ 为两个\Term{ringed}{赋环} $\mathcal U$-\Term{fibered-topos}{纤维拓扑斯}，$m:(p:F\to I,A)\to(q:G\to I,B)$ 为\Term{ringed}{赋环} $\mathcal U$-\Term{fibered-topos}{纤维拓扑斯}的\Term{morphism}{态射}。假定对 $I$ 的每个\Term{morphism}{态射} $f:i\to j$，模的\Term{functor}{函子} $f_*:\operatorname{Mod}(F_i,A_i)\to\operatorname{Mod}(F_j,A_j)$ 的\Term{derived-functor}{导出函子} $\mathrm R^nf_*$，$n\in\mathbf N$，与\Term{filtered-inductive-limit}{滤过归纳极限}可交换；对每个\Term{object}{对象} $i\in I$，模的\Term{functor}{函子} $m_{i*}:(F_i,A_i)\to(G_i,B_j)$ 的\Term{derived-functor}{导出函子} $\mathrm R^nm_{i*}$ 也与\Term{filtered-inductive-limit}{滤过归纳极限}可交换。\footnote{为明确起见，我们取左模。}回顾，当\Term{category}{范畴} $F_i,G_i$（$i\in\operatorname{ob}(I)$）\Term{coherent}{凝聚}，且\Term{morphism}{态射} $f_\bullet:F_i\to F_j$（$f:i\to j\in\operatorname{Fl}(I)$）与 $m_i:F_i\to G_i$（$i\in I$）是\Term{coherent-morphism}{凝聚态射}时，这些假设成立（5.2）。使用 8.5.1 的记号。另以 $\underrightarrow A$（相应地，$\underrightarrow B$）表示\Term{sheaf}{层} $\varinjlim\mu_i^*(A_i)$（相应地，$\varinjlim\mu_i^*(B_i)$）。最后，对于模，以 $\mu_i^{-1}$ 表示\Term{abelian-sheaf}{阿贝尔层}意义下的\Term{inverse-image}{逆像}，把 $\mu_i^*$ 专留给模意义下的\Term{inverse-image}{逆像}（第 IV 讲 11）。

\begin{SGAresult}{引理}{8.7.2}
设 $j\mapsto M_j$ 是 $\operatorname{Top}(F)$ 的一个\Term{injective-module}{内射} $A$-模。则 $Q^*(j\mapsto M_j)$ 是对 $(\underleftarrow m)_*$ \Term{acyclic}{无上同调}的 $\underrightarrow A$-模，而且对每个 $i\in\operatorname{ob}(I)$，$M_i$ 都是\Term{flasque}{松弛的}。
\end{SGAresult}
\begin{proof}
设 $i\in\operatorname{ob}(I)$，$e_i$ 是 $F_i$ 的一个\Term{final-object}{终对象}。$\mathcal U$-\Term{fibered-topos}{纤维拓扑斯} $F_{/e_i}\to I_{/i}$ 由 $F\to I$ 经\Term{base-change}{基变换} $I_{/i}\to I$ 得到。\Term{sheaf}{层} $j\mapsto M_j$ 是\Term{injective-module}{内射的}，因而是\Term{flasque}{松弛的}（第 V 讲 4.6）。它\Term{restriction}{限制}到\Term{localized-topos}{局部化拓扑斯} $F_{/e_i}$ 上仍\Term{flasque}{松弛}（第 V 讲 4.12）。从 $F_{/e_i}$ 到 $F_i$ 的\Term{restriction}{限制}\Term{functor}{函子}是一个\Term{topos}{拓扑斯}\Term{morphism}{态射}的\Term{direct-image-functor}{正像函子}（7.4.12），因此把\Term{flasque}{松弛模}变成\Term{flasque}{松弛模}（第 V 讲 5.2）。故 $M_i$ 是\Term{flasque}{松弛的}。

证明引理的第一个断言。置 $N'=Q^*(j\mapsto M_j)$。由（8.5.3.1），对每个 $i\in\operatorname{ob}(I)$，有
\[
 \mu_{i*}(N')=\mu_{i*}Q^*(j\mapsto M_j)\simeq\varinjlim_{f:j\to i}f_*(M_j).
\]
利用假设（8.7.1），可见\Term{sheaf}{层} $N'$ 具有下列性质：

（P）对每个\Term{object}{对象} $i\in\operatorname{ob}(I)$，$\mu_{i*}(N')$ 对 $m_{i*}$ \Term{acyclic}{无上同调}；对 $I$ 的每个\Term{morphism}{态射} $f:i\to j$，$\mu_{i*}(N')$ 对 $f_*$ \Term{acyclic}{无上同调}。

只须证明每个具有性质（P）的\Term{object}{对象} $N'$ 都对 $(\underleftarrow m)_*$ \Term{acyclic}{无上同调}。由于\Term{injective-module}{内射对象}具有性质（P），只须核实以下两个性质（第 V 讲 0.4）：
\SourcePage{230}
对每个\Term{exact-sequence}{正合列} $0\to N'\to N''\to0$，其中 $N'$ 与 $N$ 具有性质（P），（a）$N''$ 具有性质（P）；（b）$m_*(N)\to m_*(N'')$ 是\Term{epimorphism}{满态射}。（a）的核实是平凡的。现核实（b）。置 $K=\operatorname{coker}(Q_*(N')\to Q_*(N))$，于是 $K$ 为\Term{section}{截面} $(i\mapsto K_i=\operatorname{coker}(\mu_{i*}(N')\to\mu_{i*}(N)))$。对每个 $f:i\to j$，有\Term{commutative-diagram}{交换图}
\[
\begin{tikzcd}[column sep=small]
0\arrow[r]&f_*\mu_{i*}(N')\arrow[r]&f_*\mu_{i*}(N)\arrow[r]&f_*(K_i)\arrow[r]&0\\
0\arrow[r]&\mu_{j*}(N')\arrow[r]\arrow[u]&\mu_{j*}(N)\arrow[r]\arrow[u]&K_j\arrow[r]\arrow[u]&0.
\end{tikzcd}
\]
因 $\mu_{i*}(N')$ 对 $f_*$ \Term{acyclic}{无上同调}，上方横列\Term{exact}{正合}。因 $i\mapsto\mu_{i*}(N')$ 与 $i\mapsto\mu_{i*}(N)$ 是\Term{cartesian-section}{笛卡尔截面}，前两个竖\Term{morphism}{态射}是\Term{isomorphism}{同构}。因此\Term{canonical}{典范}\Term{morphism}{态射} $K_j\to f_*(K_i)$ 是\Term{isomorphism}{同构}，$i\mapsto K_i$ 是\Term{cartesian-section}{笛卡尔截面}。故 $K$ 形如 $Q_*(K')$；又因 $Q_*$ \Term{fully-faithful}{全忠实}，\Term{canonical}{典范}\Term{morphism}{态射} $Q_*(N)\to K$ 与 $K\to Q_*(N'')$ 分别形如 $Q_*(u)$ 和 $Q_*(v)$。因 $Q^*$ \Term{exact}{正合}，且 $Q^*Q_*$ \Term{isomorphism}{同构}于\Term{identity-functor}{恒等函子}，$v$ 是\Term{isomorphism}{同构}，从而 $Q_*(v)$ 也是\Term{isomorphism}{同构}。因此列 $0\to Q_*(N')\to Q_*(N)\to Q_*(N'')\to0$ \Term{exact}{正合}，进而对 $I$ 的每个\Term{object}{对象} $i$，列 $0\to\mu_{i*}(N)\to\mu_{i*}(N')\to\mu_{i*}(N'')\to0$ \Term{exact}{正合}。由于 $\mu_{i*}(N')$ 对 $m_{i*}$ \Term{acyclic}{无上同调}，列
\[
 0\to m_{i*}\mu_{i*}(N')\to m_{i*}\mu_{i*}(N)\to m_{i*}\mu_{i*}(N'')\to0
\]
\Term{exact}{正合}。又有\Term{isomorphism}{同构} $(\underleftarrow m)_*\simeq\varinjlim_{I^\circ}\mu_j^{-1}m_{j*}\mu_{j*}$（8.5.8.2）。\Term{functor}{函子} $\mu_j^{-1}$ \Term{exact}{正合}，而\Term{filtered-inductive-limit}{滤过归纳极限}是\Term{exact}{正合的}。故列 $0\to(\underleftarrow m)_*(N')\to(\underleftarrow m)_*(N)\to(\underleftarrow m)_*(N'')\to0$ \Term{exact}{正合}。
\end{proof}

\begin{SGAresult}{定理}{8.7.3}
采用 8.7.1 的记号与假设。对每个整数 $n\in\mathbf N$，\Term{derived-functor}{导出函子} $\mathrm R^n(\underleftarrow m)_*$ 与\Term{filtered-inductive-limit}{滤过归纳极限}可交换；对每个 $A$-模 $j\mapsto M_j$，有\Term{functorial}{函子性}\Term{isomorphism}{同构}
\[
 \mathrm R^n(\underleftarrow m)_*Q^*(j\longmapsto M_j)
 \simeq\varinjlim_{I^\circ}\mu_j^*\mathrm R^nm_{j*}(M_j).\tag{8.7.3.1}
\]
\end{SGAresult}
\begin{proof}
有 $Q^{-1}(A)=\underrightarrow A$，所以模的\Term{inverse-image}{逆像}\Term{isomorphism}{同构}于\Term{abelian-sheaf}{阿贝尔层}的\Term{inverse-image}{逆像}。由 8.5.8.2 可得，对每个\Term{section}{截面} $j\mapsto M_j$，有\Term{canonical-isomorphism}{典范同构} $\varinjlim_{I^\circ}\mu_j^{-1}(M_j)\simeq\varinjlim_{I^\circ}\mu_j^*(M_j)$。又在 8.7.1 的假设下，有\Term{canonical-isomorphism}{典范同构}（8.5.9.1）
\[
 (\underleftarrow m)_*Q^*(j\mapsto M_j)\simeq
 \varinjlim_{I^\circ}\mu_j^{-1}m_{j*}(M_j).
\]
由于对每个\Term{injective-module}{内射对象} $j\mapsto M_j$，各 $M_j$ 都\Term{flasque}{松弛}，且 $Q^*(j\mapsto M_j)$ 对 $(\underleftarrow m)_*$ \Term{acyclic}{无上同调}（8.7.2），有\Term{isomorphism}{同构}
\[
 \mathrm R^n(\underleftarrow m)_*(Q^*(j\mapsto M_j))\simeq
 \varinjlim_{I^\circ}\mu_j^{-1}\mathrm R^nm_{j*}(M_j).
\]
再利用上述事实即得（8.7.3.1）。第一个断言直接由公式（8.7.3.1）得出。
\end{proof}

\begin{SGAresult}{推论}{8.7.4}
对 $n\in\mathbf Z$，有\Term{functorial}{函子性}\Term{isomorphism}{同构}
\[
 \mathrm R^n(\underleftarrow m)_*(N)\simeq
 \varinjlim_{I^\circ}\mu_j^*\mathrm R^nm_{j*}\mu_{j*}(N).\tag{8.7.4.1}
\]
\end{SGAresult}
\SourcePage{231}
\begin{proof}
因为 $N$ \Term{isomorphism}{同构}于 $Q^*(j\mapsto\mu_{j*}(N))$。
\end{proof}

\begin{SGAresult}{推论}{8.7.5}
对 $I$ 的每个\Term{object}{对象} $i$、每个整数 $n$ 及每个 $A_i$-模 $M_i$，有\Term{functorial}{函子性}\Term{isomorphism}{同构}
\[
 \mathrm R^n(\underleftarrow m)_*\mu_i^*(M_i)\simeq
 \varinjlim_{f:j\to i}\mu_j^*\mathrm R^nm_{j*}f^*(M_i).\tag{8.7.5.1}
\]
\end{SGAresult}
\begin{proof}
作\Term{base-change}{基变换} $I_{/i}\to I$，可假定 $i$ 为 $I$ 的\Term{final-object}{终对象}（8.2.1）。此时\Term{object}{对象} $\mu_i^*(M_i)$ \Term{isomorphism}{同构}于 $Q^*((f:j\to i)\mapsto f^*(M_i))$（参见 8.5.5）。
\end{proof}

\begin{SGAresult}{推论}{8.7.6}
设 $i$ 为 $I$ 的\Term{object}{对象}，$n$ 为整数。关于\Term{section}{截面} $j\mapsto M_j$，有\Term{functorial}{函子性}\Term{isomorphism}{同构}
\[
 \mathrm R^nm_{i*}Q^*(j\mapsto M_j)\simeq
 \varinjlim_{f:j\to i}\mathrm R^nf_*(M_j).\tag{8.7.6.1}
\]
关于 $\underrightarrow A$-模 $N$，有\Term{isomorphism}{同构}
\[
 \mathrm R^n\mu_{i*}(N)\simeq\varinjlim_{f:j\to i}\mathrm R^nf_*\mu_{j*}(N).\tag{8.7.6.2}
\]
关于 $A_i$-模 $M_i$，有\Term{functorial}{函子性}\Term{isomorphism}{同构}
\[
 \mathrm R^n\mu_{i*}\mu_i^*(M_i)\simeq
 \varinjlim_{f:j\to i}\mathrm R^nf_*f_*(M_i).\tag{8.7.6.3}
\]
\Term{functor}{函子} $\mathrm R^n\mu_{i*}$ 与\Term{filtered-inductive-limit}{滤过归纳极限}可交换。
\end{SGAresult}
\begin{proof}
作\Term{base-change}{基变换} $I_{/i}\to I$，可归结到 $i$ 为 $I$ 的\Term{final-object}{终对象}的情形。取 $G$ 为\Term{fiber}{纤维} $F_i$ 的常值\Term{fibered-topos}{纤维拓扑斯}，取\Term{morphism}{态射} $m$ 为 $(f_\bullet,f\in\operatorname{ob}(I_{/i}))$，则公式（8.7.6.1）和（8.7.6.2）分别是公式（8.7.3.1）、（8.7.4.1）和（8.7.5.1）的特殊情形。最后一个断言由（8.7.6.2）并结合 8.5.4 得出。
\end{proof}

\begin{SGAresult}{推论}{8.7.7}
设 $i$ 为 $I$ 的\Term{object}{对象}，$X$ 为 $F_i$ 的\Term{object}{对象}，且\Term{functor}{函子} $\mathrm H^n(X,-)$，$n\in\mathbf N$，与\Term{filtered-inductive-limit}{滤过归纳极限}可交换。对每个 $n$，关于 $A_i$-模 $M_i$ 有\Term{canonical-functor}{典范函子}性\Term{isomorphism}{同构}
\[
 \mathrm H^n(\mu_i^*(X),\mu_i^*(M_i))\simeq
 \varinjlim_{f:j\to i}\mathrm H^n(f^*(X),f^*(M_i)),\tag{8.7.7.1}
\]
且\Term{functor}{函子} $\mathrm H^n(\mu_i^*(X),-)$ 与\Term{filtered-inductive-limit}{滤过归纳极限}可交换。特别地，若\Term{functor}{函子} $\mathrm H^n(F_i,-)$ 与\Term{filtered-inductive-limit}{滤过归纳极限}可交换，则\Term{functor}{函子} $\mathrm H^n(\underleftarrow F,-)$ 与\Term{filtered-inductive-limit}{滤过归纳极限}可交换，并且关于 $A_i$-模 $M_i$ 有\Term{canonical-functor}{典范函子}性\Term{isomorphism}{同构}
\[
 \mathrm H^n(\underleftarrow F,\mu_i^*(M_i))\simeq
 \varinjlim_{f:j\to i}\mathrm H^n(F_j,f^*(M_i)).\tag{8.7.7.2}
\]
\end{SGAresult}
\SourcePage{232}
\begin{proof}
由第 V 讲 5.3，有两个\Term{spectral-sequence}{谱序列}
\[
 {}'E_2^{p,q}=\varinjlim_{f:j\to i}\mathrm H^p(X,\mathrm R^qf_*f^*(M_i))
 \Longrightarrow\varinjlim_{f:j\to i}\mathrm H^{p+q}(f^*(X),f^*(M_i)),
\]
\[
 {}''E_2^{p,q}=\mathrm H^p(X,\mathrm R^q\mu_{i*}\mu_i^*(M_i))
 \Longrightarrow\mathrm H^{p+q}(\mu_i^*(X),\mu_i^*(M_i));
\]
由（8.5.7.1），有这两个\Term{spectral-sequence}{谱序列}之间的一个\Term{morphism}{态射}。由于 $\mathrm H^{p+q}(X,-)$ 与\Term{filtered-inductive-limit}{滤过归纳极限}可交换，由（8.7.6.2）得，此\Term{spectral-sequence}{谱序列}\Term{morphism}{态射}在 $E_2^{p,q}$ 页是\Term{isomorphism}{同构}。因此它在\Term{abutment}{终项}上诱导\Term{isomorphism}{同构}，从而得到（8.7.7.1）。对每个 $\underrightarrow A$-模 $N$，有\Term{spectral-sequence}{谱序列}（第 V 讲 5.3）
\[
 E_2^{p,q}=\mathrm H^p(X,\mathrm R^q\mu_{i*}(N))
 \Longrightarrow\mathrm H^{p+q}(\mu_i^*(X),N).
\]
\Term{functor}{函子} $\mathrm H^n(\mu_i^*(X),-)$ 与\Term{filtered-inductive-limit}{滤过归纳极限}的交换性，遂由\Term{functor}{函子} $\mathrm H^n(X,-)$ 及 $\mathrm R^n\mu_{i*}$ 的相应性质得出（8.7.6）。
\end{proof}

\textbf{注记 8.7.8。} 当各\Term{topos}{拓扑斯} $F_j$，$j\in\operatorname{ob}(I)$，\Term{coherent}{凝聚}，且\Term{transition-morphism}{转移态射}是\Term{coherent-morphism}{凝聚态射}时，8.6.7 中对 $X$（相应地，$F_i$）作的假设在 $X$（相应地，$F_i$）\Term{coherent}{凝聚}时成立（5.2）。另一方面，已知此时 $\mu_i^*(X)$（相应地，$\underleftarrow F$）\Term{coherent}{凝聚}（8.3.13）。

\begin{SGAresult}{推论}{8.7.9}
设 $i$ 为 $I$ 的\Term{object}{对象}，$M_i$ 为左 $A_i$-模，$L_i$ 为右 $A_i$-模，并有如下类型的\Term{resolution}{解消}：
\[
 P_{i,k}\longrightarrow P_{i,k-1}\longrightarrow\cdots\longrightarrow P_{i,0}
 \longrightarrow L_i\longrightarrow0,
\]
其中对每个整数 $k$，$P_{i,k}$ \Term{isomorphism}{同构}于形如 $A_{i\mid X}$（第 IV 讲 11.3.3）的\Term{object}{对象}的\Term{finite}{有限}\Term{direct-sum}{直和}，这里 $X$ 满足 8.7.7 的假设。若除 8.7.1 的假设外，\Term{fibered-topos}{纤维拓扑斯} $F$ 还\Term{right-flat}{右平坦}（8.6），则对 $n\leq k-1$，有\Term{canonical-isomorphism}{典范同构}
\[
 \operatorname{Ext}_{\underrightarrow A}^n(\underleftarrow F,\mu_i^*(L_i),\mu_i^*(M_i))
 \simeq\varinjlim_{f:j\to i}\operatorname{Ext}_{A_j}^n(F_j,f^*(L_i),f^*(M_i)).
 \tag{8.7.9.1}
\]
\end{SGAresult}
\begin{proof}
记 $P_{i,\bullet}$ 为 $L_i$ 的\Term{resolution}{解消}。因 $F$ \Term{right-flat}{右平坦}，对每个 $f:j\to i$，$f^*(P_{i,\bullet})$ 是 $f^*(L_i)$ 的\Term{resolution}{解消}，而 $\mu_i^*(P_{i,\bullet})$ 是 $\mu_i^*(L_i)$ 的\Term{resolution}{解消}（8.6）。这些\Term{resolution}{解消}可以构造两个\Term{spectral-sequence}{谱序列}，分别收敛到（8.7.9.1）的两端，并构造这两个\Term{spectral-sequence}{谱序列}之间的一个\Term{morphism}{态射}。此\Term{morphism}{态射}在 $E_1^{p,q}$ 页是\Term{isomorphism}{同构}（8.7.7.1），所以在\Term{abutment}{终项}上诱导\Term{isomorphism}{同构}。
\end{proof}

\SourcePage{233}
\SGAsection{9}{附录：点存在的判据}
\begin{center}P. Deligne\end{center}
\begin{SGAresult}{命题}{9.0}
每个\Term{locally-coherent}{局部凝聚}\Term{topos}{拓扑斯} $S$ 都有\Term{enough-points}{足够多的点}。
\end{SGAresult}
\begin{proof}
问题在 $S$ 上是\Term{local-scheme}{局部}的，可假定\Term{topos}{拓扑斯} $S$ 由一个\Term{site}{景} $\mathscr S$ 定义，其中\Term{finite}{有限}\Term{projective-limit}{投射极限}\Term{representable}{可表}，且每个\Term{covering}{覆盖} $f_i:U_i\to U$ 都有一个\Term{finite}{有限}\Term{subcovering}{子覆盖}。只须证明，若 $f:\mathscr F\to\mathscr G$ 不是\Term{monomorphism}{单态射}，则存在 $\mathscr S$ 的一个点 $x$，使 $f_x$ 不为\Term{injectivity}{单射}。由假设，存在 $U\in\operatorname{ob}\mathscr S$ 与 $s,s'\in\mathscr F(U)$，使 $s\ne s'$ 且 $f(s)=f(s')$。把 $\mathscr S$ 换成 $\mathscr S_{/U}$，便归结到下面的引理。
\end{proof}

\begin{SGAresult}{引理}{9.1}
若 $s,s'$ 为满足上述假设的\Term{site}{景} $\mathscr S$ 上的\Term{sheaf}{层} $\mathscr F$ 的两个不同的\Term{global-section}{整体截面}，则存在 $\mathscr S$ 的一个点 $x$，使 $s_x\ne s'_x$。
\end{SGAresult}
\begin{proof}
设 $P=(U_i)_{i\in I}$ 为 $\mathscr S$ 中的一个\Term{projective-system}{投射系统}，由\Term{filtered-ordered-set}{滤过有序集} $I$ 编号。对 $\mathscr S$ 上的每个\Term{sheaf}{层} $\mathscr H$，置 $P(\mathscr H)=\varinjlim\mathscr H(U_i)$；对每个 $V\in\operatorname{Ob}\mathscr S$，置 $P(V)=\varinjlim\operatorname{Hom}(U_i,V)$。\Term{functor}{函子} $P(\mathscr H)$ 与 $P(V)$ 均与\Term{fiber-product}{纤维积}可交换。若\Term{functor}{函子} $P(V)$ 把\Term{covering}{覆盖}变成映射的\Term{surjective-family}{满射族}，则它是从\Term{point-topos}{点拓扑斯}到 $\mathscr S$ 的\Term{site-morphism}{景态射}，而 $P(\mathscr H)$ 就是对应的\Term{fiber-functor}{纤维函子}。

若 $P=(U_i)_{i\in I}$ 与 $Q=(V_j)_{j\in J}$ 是 $\mathscr S$ 中的两个\Term{projective-system}{投射系统}，当 $I$ 是 $J$ 的一个\Term{subset}{子集}、配备诱导的序，且 $P$ 为 $Q$ 对 $J$ 的\Term{restriction}{限制}时，称 $Q$ \Term{refinement}{加细} $P$。此时有从 $P(\mathscr H)$ 到 $Q(\mathscr H)$ 及从 $P(V)$ 到 $Q(V)$ 的\Term{morphism-of-functors}{函子态射}。

对上述任意 $P$，以 $s_P,s'_P$ 表示 $s,s'$ 在 $P(\mathscr F)$ 中的像。在下面的引理中，取 $P$ 为由 $\{0\}$ 编号且只含 $\mathscr S$ 的\Term{final-object}{终对象}的\Term{projective-system}{投射系统}，即可得到引理 9.1。
\end{proof}

\begin{SGAresult}{引理}{9.2}
对上述每个满足 $s_P\ne s'_P$ 的 $P$，存在\Term{refinement}{加细} $P$ 的 $Q$，满足 $s_Q\ne s'_P$，并且对任意\Term{covering}{覆盖} $f_i:V_i\to V$ 与任意 $f\in Q(V)$，$f$ 都属于某个 $Q(V_i)$ 的像。
\end{SGAresult}
先证明：
\begin{SGAresult}{引理}{9.3}
设 $P=(U_i)_{i\in I}$ 满足引理 9.2 的假设，$f_i:V_i\to V$ 是 $\mathscr S$ 中的\Term{finite}{有限}\Term{covering}{覆盖}，$f\in P(V)$。存在\Term{refinement}{加细} $P$ 的 $Q$，满足 $s_Q\ne s'_Q$，使得 $f$ 在 $Q(V)$ 中的像属于某个 $(V_i)$ 的像。
\end{SGAresult}
\begin{proof}
存在 $i_\circ\in I$ 和 $f'\in\operatorname{Hom}(U_{i_\circ},V)$ 定义 $f$。对 $i\geq i_\circ$，置 $U_{i,k}=U_i\times_VV_k$，此\Term{fiber-product}{纤维积}借 $f'$ 定义。已知 $(V_k)$ 是 $V$ 的\Term{covering}{覆盖}，
\SourcePage{234}
所以 $(U_{i,k})_k$ 是 $U_i$ 的\Term{covering}{覆盖}；以下\Term{arrow}{箭头}将是\Term{injectivity}{单射}（$i>i_\circ$）：
\[
 \mathscr F(U_i)\longrightarrow\prod_k\mathscr F(U_{i,k}).
\]
取\Term{limit}{极限}，并利用\Term{finite-product}{有限积}与\Term{filtered-inductive-limit}{滤过归纳极限}可交换，可见
\[
 P(\mathscr F)\longrightarrow\prod_k\varinjlim\mathscr F(U_{i,k})
\]
是\Term{injectivity}{单射}，因而存在 $k$，使 $s_P,s'_P$ 在 $\varinjlim_{i>i_\circ}\mathscr F(U_{i,k})$ 中的像不同。置 $I_\circ=\{i\mid i\geq i_\circ\text{ 且 }i\in I\}$，$J=I\amalg I_\circ$。在 $J$ 上定义序：$j'\leq j''$ 当且仅当 $j'$ 与 $j''$ 在 $I$ 中的像满足 $j'\leq j''$，且并非 $j'\in I_\circ$、$j''\in I$。由 $J$ 编号的 $U_i$ 与 $U_{i,k}$ 构成 $P$ 的一个\Term{refinement}{加细} $Q$，满足 $s_Q\ne s'_Q$，且 $f$ 在 $Q(V)$ 中的像属于 $Q(V_k)$ 的像。
\end{proof}

\begin{SGAresult}{引理}{9.4}
设 $P=(U_i)_{i\in I}$ 满足引理 9.2 的假设。存在\Term{refinement}{加细} $P$ 的 $Q$，满足 $s_Q\ne s'_Q$，使对于 $S$ 的任意\Term{open-subset}{开集} $V$ 的任意\Term{finite}{有限}\Term{covering}{覆盖} $(V_k)$ 及任意 $f\in P(V)$，$f$ 在 $Q(V)$ 中的像属于某个 $Q(V_k)$ 的像。
\end{SGAresult}
\begin{proof}
设 $E$ 为由 $S$ 的一个\Term{open-subset}{开集} $V$、$V$ 的一个\Term{finite}{有限}\Term{covering}{覆盖} $(V_k)$ 及 $f\in P(V)$ 组成的三元组的\Term{set}{集合}。设 $\leq$ 为 $E$ 上的一个\Term{well-order}{良序}，$\overline E$ 为在 $E$ 中添入最大元素 $\infty$ 所得的\Term{set}{集合}。对 $e\in\overline E$ 作\Term{transfinite-induction}{超限归纳}，定义 $P$ 的\Term{refinement}{加细} $Q_e$，满足 $s_{Q_e}\ne s'_{Q_e}$，且：

（i）若 $e'\leq e$，则 $Q_e$ \Term{refinement}{加细} $Q_{e'}$。

（ii）若 $e=(V,(V_k),f)$，则 $f$ 在 $Q_e(V)$ 中的像属于某个 $Q_e(V_k)$ 的像。

假定对于 $e'<e$，$Q_{e'}$ 已经定义。若 $e$ 为 $E$ 的首元素，置 $Q'_e=P$。若 $e$ 有前驱 $e-1$，置 $Q'_e=Q_{e-1}$。否则，取 $Q'_e$ 为由指标集 $\bigcup_{e'<e}I_{e'}$ 编号、\Term{refinement}{加细}所有 $e'<e$ 的 $Q_{e'}$ 的\Term{projective-system}{投射系统}。此时有
\[
 Q'_e(\mathscr F)=\varinjlim_{e'<e}Q_{e'}(\mathscr F),
\]
因而无论哪种情形，$Q'_e$ 都\Term{refinement}{加细}所有 $e'<e$ 的 $Q_{e'}$，且满足 $s_{Q'_e}\ne s'_{Q'_e}$。

对 $Q'_e$ 与 $e=(V,(V_k),f)$ 应用引理 9.3，就得到所需的\Term{projective-system}{投射系统} $Q_e$（若 $e=\infty$，则取 $Q_e=Q'_e$）。

\Term{projective-system}{投射系统} $Q_\infty$ 满足引理 9.4。
\end{proof}
\begin{proof}[引理 9.2 的证明]
引理 9.4 允许对 $n$ 归纳定义 $\mathscr S$ 中的\Term{projective-system}{投射系统}序列 $Q_n=(U_i)_{i\in I_n}$，使：

（i）$Q_0=P$；

（ii）$Q_{n+1}$ \Term{refinement}{加细} $Q_n$；

（iii）$s_{Q_n}\ne s'_{Q_n}$；
\SourcePage{235}

（iv）对任意\Term{covering}{覆盖} $f_k:V_k\to V$ 与 $f\in Q_n(V)$，$f$ 在 $Q_{n+1}(V)$ 中的像属于某个 $Q_{n+1}(V_k)$ 的像。

以所有 $I_n$ 的并为指标集、延拓各个 $Q_n$ 的\Term{projective-system}{投射系统} $Q$，便满足引理 9.2。证明还表明：
\end{proof}

\begin{SGAresult}{推论}{9.5}
设 $\mathscr S$ 为一个\Term{site}{景}，其中\Term{fiber-product}{纤维积}\Term{representable}{可表}，且 $\mathscr S$ 中的每个\Term{covering}{覆盖}都有\Term{finite}{有限}\Term{subcovering}{子覆盖}。设 $e$ 为\Term{cardinal}{基数}
\[
 \sup\bigl((\text{\SGAsourcegap{原页此处基数符号缺失}})_0,
 \operatorname{card}(\operatorname{Fl}\mathscr S)\bigr).
\]
则存在 $\mathscr S$ 的一个\Term{very-dense-points}{极稠密点集} $X$，使得
\[
 \text{（i）}\quad\operatorname{card}(X)\leq2^c;\qquad
 \text{（ii）}\quad x\in X,\ U\in\operatorname{Ob}\mathscr S
 \ \Longrightarrow\ \operatorname{card}(U_x)\leq c.
\]
\end{SGAresult}

\begin{center}\textbf{参考文献}\end{center}
\begin{enumerate}
\item P. Gabriel 与 M. Zisman：《\Term{homotopy}{同伦}论与分式演算》（Homotopy Theory and Calculus of Fraction），Ergebnisse der Mathematik，第 35 卷。
\item J. Giraud：《\Term{descent}{下降}方法》（Méthode de la descente），Mémoire de la S.M.F.
\item A. Grothendieck 与 J. Dieudonné：《\Term{algebra}{代数}几何基础》I（Éléments de Géométrie Algébrique I），第 2 版。
\item A. Grothendieck 与 J. Dieudonné：《\Term{algebra}{代数}几何基础》IV I，I.H.E.S. n\textsuperscript{o}（原文未填期号）。
\item SGA 1，第 VI 讲，A. Grothendieck，I.H.E.S.
\item SGA 3，第 IV 讲 6.3，M. Demazure，Lecture Notes n\textsuperscript{o}（原文未填卷号），Springer Verlag。
\item A. Grothendieck：《关于\Term{homology}{同调}\Term{algebra}{代数}的若干问题》（Sur quelques points d'algèbres homologiques），Tohoku Math. J.
\end{enumerate}
\SourcePage{236}
% 原 PDF 此页空白。

\SourcePage{237}
\SGAchapter{VII}{概形的平展景与平展拓扑斯}{Site et topos étales d'un schéma}
\begin{center}A. Grothendieck\end{center}
% 本页原目录由共享模板按中文标题和译本页码重建。
本讲和下一讲将阐述\Term{etale-topology}{平展拓扑}与\Term{etale-cohomology}{平展上同调}最基本的性质。本讲讨论的若干性质，本质上也适用于其他十分不同的\Term{topology}{拓扑}，例如“fppf \Term{topology}{拓扑}”。下一讲将讨论\Term{etale-topology}{平展拓扑}特有的一些性质，它们来自\Term{etale-morphism}{平展态射}十分特殊的本性。

这里部分遵循 J. E. Roos 的三次口头报告（他本人未能将其写成文字），尤其是第 VIII 讲 6.3 的证明。

除非明确另有说明，本讲及以后各讲所考虑的\Term{scheme}{概形}均默认为固定\Term{universe}{宇宙} $\mathcal U$ 的元素。

\SGAsection{1}{平展拓扑}
\textbf{1.1.} 以 $(\mathrm{Sch})$ 表示\Term{scheme}{概形}（固定\Term{universe}{宇宙} $\mathcal U$ 的元素）的\Term{category}{范畴}。回顾，$(\mathrm{Sch})$ 中的\Term{morphism}{态射} $f:X\to Y$ 称为\Term{etale}{平展的}，若它\Term{locally-finite-presentation}{局部有限表现}、\Term{flat}{平坦}，而且对每个 $y\in Y$，\Term{fiber}{纤维} $X_y$ 是\Term{discrete}{离散}的，其\Term{local-ring}{局部环}都是 $k(y)$ 的\Term{finite}{有限}\Term{separable-extension}{可分扩张}。\Term{equivalence}{等价}地，$f$
\SourcePage{238}
\Term{locally-finite-presentation}{局部有限表现}，并且对每个（绝对意义下）\Term{affine}{仿射}的 $Y$-\Term{scheme}{概形} $Y'$，以及每个由\Term{nilpotent-ideal}{幂零理想}定义的\Term{subscheme}{子概形} $Y'_0$，映射
\[
 \operatorname{Hom}_Y(Y',X)\longrightarrow\operatorname{Hom}_Y(Y'_0,X)
\]
是\Term{bijection}{双射}。关于这一概念最重要的性质，参见 EGA IV 第 17、18 节；在其出版之前，可参见 SGA 1 第 I、IV 讲。

\textbf{1.2.} $(\mathrm{Sch})$ 上的\Term{etale-topology}{平展拓扑}，是由下列\Term{pretopology}{预拓扑}生成的\Term{topology}{拓扑}：对每个 $X\in(\mathrm{Sch})$，\Term{set}{集合} $\operatorname{Cov}X$ 由所有族 $(X_i\xrightarrow{u_i}X)_{i\in I}$（由 $I\in\mathcal U$ 编号）组成，其中各 $u_i$ \Term{etale}{平展}，且在\Term{set}{集合}意义下 $X=\bigcup_i u_i(X_i)$。

给每个 $X$ 配上 $(\mathrm{Sch})_{/X}$ 中由\Term{etale}{平展}\Term{arrow}{箭头} $X'\to X$ 组成的\Term{subcategory}{子范畴} $\mathrm{Et}_{/X}$，是很方便的。赋予它由 $(\mathrm{Sch})$ 的\Term{etale-topology}{平展拓扑}“诱导的\Term{topology}{拓扑}”（第 III 讲），也称为“$X$ 上的\Term{etale-topology}{平展拓扑}”；所得\Term{site}{景}记为 $X_{\text{ét}}$（$X$ 的“\Term{etale-site}{平展景}”）。\footnote{不过，最好用 $X_{\text{ét}}$ 表示由 $X$ 的\Term{etale-site}{平展景}定义的\Term{topos}{拓扑斯} $\widetilde X_{\text{ét}}$。出于实用考虑，本研讨班仍沿用这里引入的记号，即原研讨班的记号。}注意 $\mathrm{Et}_{/X}$ 的每个\Term{morphism}{态射}都\Term{etale}{平展}（SGA 1 第 I 讲 4.8），所以 $X_{\text{ét}}$ 中的族 $(X'_i\xrightarrow{u_i}X')_{i\in I}$ 是\Term{covering-family}{覆盖族}，当且仅当它是\Term{surjective-family}{满射族}，即 $\bigcup_i u_i(X'_i)=X'$。立即可见（参见 3.1），$X_{\text{ét}}$ 是 $\mathcal U$-\Term{site}{景}，故第 I 至 VI 讲的结果均可适用。$X_{\text{ét}}$ 上的 $\mathcal U$-\Term{sheaf}{层}所成的\Term{topos}{拓扑斯} $\widetilde X_{\text{ét}}$，就是 $X$ 的\Term{etale-topos}{平展拓扑斯}。

\textbf{1.3.} 以下除非明确另有说明，$\mathrm{Et}_{/X}$ 或 $(\mathrm{Sch})$ 中所考虑的一切\Term{topology}{拓扑}概念，都按\Term{etale-topology}{平展拓扑}理解。此外，语言和记号中常以 $X$ 代替 $X_{\text{ét}}$ 或 $\mathrm{Et}_{/X}$。因此，“$X$ 上的\Term{sheaf}{层}”（默认为\Term{etale-topology}{平展拓扑}）指\Term{site}{景} $X_{\text{ét}}$ 上的\Term{sheaf}{层}。以 $\widetilde X_{\text{ét}}$ 表示这些\Term{sheaf}{层}的\Term{category}{范畴}，它是一个 $\mathcal U$-\Term{topos}{拓扑斯}。若 $F$ 为 $X$ 上的\Term{abelian-sheaf}{阿贝尔层}，其\Term{cohomology-group}{上同调群}记为 $\mathrm H^i(X,F)$，在第 V 讲中它们的记号为 $\mathrm H^i(X_{\text{ét}},F)$。

\textbf{1.4.} 设 $f:X\to Y$ 为\Term{scheme}{概形}\Term{morphism}{态射}。则\Term{inverse-image-functor}{逆像函子} $f^*:Y'\mapsto Y'\times_YX$ 诱导\Term{functor}{函子}
\[
 f_{\text{ét}}^*:\mathrm{Et}_{/Y}\longrightarrow\mathrm{Et}_{/X},\tag{*}
\]
它与\Term{finite}{有限}\Term{projective-limit}{投射极限}可交换，且把\Term{covering-family}{覆盖族}变成\Term{covering-family}{覆盖族}，因而是一个\Term{site-morphism}{景态射}
\[
 f_{\text{ét}}^*:X_{\text{ét}}\longrightarrow Y_{\text{ét}},
\]
\SourcePage{239}
从而在\Term{sheaf}{层}\Term{category}{范畴}上诱导\Term{functor}{函子}
\[
 f_*^{\text{ét}}:\widetilde X_{\text{ét}}\longrightarrow\widetilde Y_{\text{ét}},
 \qquad f_*^{\text{ét}}(F)=F\circ f_{\text{ét}}^*.
\]
此外 $f_*^{\text{ét}}$ 有一个左\Term{adjoint-functor}{伴随函子}
\[
 f_{\text{ét}}^*:\widetilde Y_{\text{ét}}\longrightarrow\widetilde X_{\text{ét}},
\]
它延拓（$*$）中考虑的\Term{functor}{函子}，并与任意\Term{inductive-limit}{归纳极限}及\Term{finite}{有限}\Term{projective-limit}{投射极限}可交换，即 $f_{\text{ét}}^*$ 定义了一个\Term{topos}{拓扑斯}\Term{morphism}{态射}
\[
 f_{\text{ét}}:\widetilde X_{\text{ét}}\longrightarrow\widetilde Y_{\text{ét}}.
\]
显然，对两个\Term{scheme}{概形}\Term{morphism}{态射}的复合 $X\xrightarrow fY\xrightarrow gZ$，有\Term{canonical-isomorphism}{典范同构}
\[
 \begin{cases}
 (gf)_*^{\text{ét}}\simeq g_*^{\text{ét}}f_*^{\text{ét}},&\text{从而}\\
 (gf)_{\text{ét}}^*\simeq f_{\text{ét}}^*g_{\text{ét}}^*,&\text{即有同构}\\
 (gf)_{\text{ét}}\simeq g_{\text{ét}}f_{\text{ét}}.
 \end{cases}
\]
因此得到一个从\Term{category}{范畴} $(\mathrm{Sch})$ 到由属于 $\mathcal U'$ 的\Term{topos}{拓扑斯}组成的\Term{category}{范畴} $(\mathrm{Top})$ 的“\Term{pseudo-functor}{伪函子}”（SGA 1 第 VI 讲 8），其中 $\mathcal U'$ 是使 $\mathcal U\in\mathcal U'$ 成立的最小\Term{universe}{宇宙}；比较第 IV 讲。

\textbf{1.5. 记号。} 以下常以 $f^*,f_*$ 代替 $f_{\text{ét}}^*,f_*^{\text{ét}}$。若 $f:X\to Y$ 为\Term{scheme}{概形}\Term{morphism}{态射}，则\Term{functor}{函子} $\mathrm R^if_*^{\text{ét}}$ 简记为 $\mathrm R^if_*$。用这些记号，回顾对 $f$ 及 $X$ 上\Term{abelian-sheaf}{阿贝尔层} $F$ 的 Leray \Term{spectral-sequence}{谱序列}（第 V 讲 5）：
\[
 \mathrm H^*(X,F)\Longleftarrow E_2^{p,q}=\mathrm H^p(Y,\mathrm R^qf_*(F)).
\]
同样，若有两个\Term{morphism}{态射} $f:X\to Y$、$g:Y\to Z$，则有 Leray \Term{spectral-sequence}{谱序列}（\Term{composite-functor}{复合函子}的\Term{spectral-sequence}{谱序列}）
\[
 \mathrm R^*(gf)_*(F)\Longleftarrow E_2^{p,q}=\mathrm R^pg_*(\mathrm R^qf_*(F)).
\]

\textbf{1.6.} 当 $f:X\to Y$ 本身是\Term{etale-morphism}{平展态射}时，可把 $X$ 看成 $Y_{\text{ét}}$ 的\Term{object}{对象}，且有\Term{site}{景}的\Term{canonical-isomorphism}{典范同构}
\[
 X_{\text{ét}}\xrightarrow\sim(Y_{\text{ét}})_{/X}.
\]
借此\Term{isomorphism}{同构}，（$*$）中的\Term{functor}{函子} $f_{\text{ét}}^*:Y'\mapsto Y'\times_YX$ 等同于\Term{category}{范畴} $\mathrm{Et}_{/Y}$ 内部的\Term{base-change}{基变换}\Term{functor}{函子}。因此\Term{functor}{函子}
\[
 f^*:\widetilde Y_{\text{ét}}\longrightarrow\widetilde X_{\text{ét}}
\]
\SourcePage{240}
此时\Term{isomorphism}{同构}于“\Term{restriction}{限制}到 $Y$”的\Term{functor}{函子}：
\[
 f^*(F)\simeq F\circ i_{X/Y},
\]
其中 $F$ 是 $Y_{\text{ét}}$ 上的\Term{sheaf}{层}，而 $i_{X/Y}:X_{\text{ét}}\to Y_{\text{ét}}$ 是显然的\Term{functor}{函子}，即把 $X$ 上的\Term{etale}{平展}\Term{scheme}{概形} $u:X'\to X$，借 $fu:X'\to Y$ 看成 $Y$ 上的\Term{etale}{平展}\Term{scheme}{概形}。

\textbf{1.7. \Term{universe}{宇宙}问题。} 注意，若 $X\ne\varnothing$，则 $X_{\text{ét}}$ 不是所选\Term{universe}{宇宙} $\mathcal U$ 的元素。不过，如前所述，容易看出（3.1），存在 $X_{\text{ét}}$ 的一个\Term{full-subcategory}{全子范畴} $C$，它是 $\mathcal U$ 的元素，且满足“\Term{comparison-lemma}{比较引理}”（第 III 讲）的条件，使\Term{restriction}{限制}\Term{functor}{函子}诱导一个\Term{equivalence}{等价} $\widetilde C\to\widetilde X_{\text{ét}}$。因此\Term{etale-topos}{平展拓扑斯}\Term{equivalence}{等价}于由某个\Term{site}{景} $C\in\mathcal U$ 定义的\Term{topos}{拓扑斯}。换言之，$X_{\text{ét}}$ 是一个 $\mathcal U$-\Term{site}{景}，故由该处结果，第 I 至 VI 讲的结论可用于这个\Term{site}{景}。

\SGAsection{2}{层的例子}
a）设 $F\in\operatorname{Ob}(\mathrm{Sch})_{/X}$ 是 $X$ 上的\Term{scheme}{概形}；对 $X$ 上每个 $X'$，置
\[
 F(X')=\operatorname{Hom}_X(X',F).
\]
这样定义的\Term{functor}{函子} $(\mathrm{Sch})_{/X}^\circ\to(\mathrm{Ens})$ 是\Term{etale-topology}{平展拓扑}的\Term{sheaf}{层}（甚至是 SGA 3 第 IV 讲所研究的更细的 \Term{fpqc-topology}{fpqc 拓扑}的\Term{sheaf}{层}）。换言之，$(\mathrm{Sch})$ 上的\Term{etale-topology}{平展拓扑}比\Term{canonical-topology}{典范拓扑}粗。事实上，这本质上就是 SGA 1 第 VIII 讲 5.1 的内容（另见 SGA 3 第 IV 讲 6.3.1）。更有甚者，此\Term{sheaf}{层}在 $X_{\text{ét}}$ 上的\Term{restriction}{限制}也是\Term{sheaf}{层}。不致混淆时仍记作 $F$。\footnote{但须注意，若 $X\ne\varnothing$，把 $F\in\operatorname{Ob}(\mathrm{Sch})_{/X}$ 送到相应\Term{sheaf}{层} $\varphi(F)$ 的\Term{functor}{函子} $\varphi$ 并不\Term{fully-faithful}{全忠实}，甚至不\Term{faithful}{忠实}，而且不能由 $\varphi(F)$ 重建 $F$（在\Term{isomorphism}{同构}意义下）。例如，若 $S=\operatorname{Spec}k$，$k$ 为\Term{algebraically-closed-field}{代数闭域}，则知道 $\varphi(F)$ 只相当于知道 $F$ 的底层\Term{set}{集合}（由第 VIII 讲 2.4）。比较第 IV 讲。}还须注意，由此得到的\Term{functor}{函子}
\[
 (\mathrm{Sch})_{/X}\longrightarrow\widetilde X_{\text{ét}}\tag{*}
\]
与\Term{finite}{有限}\Term{projective-limit}{投射极限}可交换（这显然）。例如，这意味着，若 $F$ 是 $X$ 上的\Term{group-scheme}{群概形}（相应地，……），则它定义的\Term{sheaf}{层}是\Term{sheaf-of-groups}{群层}（相应地，……）。由（$*$）诱导的\Term{functor}{函子} $\mathrm{Et}_{/X}\to\widetilde X_{\text{ét}}$ 正是\Term{canonical-functor}{典范函子}，把每个 $X'\in\operatorname{Ob}X_{\text{ét}}$ 送到它在 $X_{\text{ét}}$ 上所\Term{representable}{表示}的\Term{functor}{函子}。因此它给出从\Term{category}{范畴} $\mathrm{Et}_{/X}$ 到\Term{etale-topos}{平展拓扑斯} $\widetilde X_{\text{ét}}$ 的一个\Term{full-subcategory}{全子范畴}的\Term{isomorphism}{同构}，我们通常借此把 $X\in\operatorname{Ob}X_{\text{ét}}$ 与相应\Term{sheaf}{层}等同，记为 $\widetilde X$ 或简记为 $X$。

显然有 $\mathrm H^0(X,F)=\operatorname{Hom}_{(\mathrm{Sch})_{/X}}(X,F)$。

当 $F$ 是 $X$ 上的\Term{group-prescheme}{群预概形}时，也\Term{solvable}{可解}释 $\mathrm H^1(X,F)$（如有必要可假定交换，使 $\mathrm H^1(X,F)$ 的定义
\SourcePage{241}
属于第 V 讲的范围；一般情形可参阅 J. Giraud 的学位论文\footnote{原书此脚注为空。}）。熟知的论证表明，我不在这里重复，$\mathrm H^1(X,F)$ \Term{canonical-isomorphism}{典范同构}于 $X_{\text{ét}}$ 上在 $F$ \Term{action}{作用}下\Term{principal-homogeneous}{主齐性}的\Term{sheaf-of-sets}{集合层}的类（模\Term{isomorphism}{同构}）组成的群。\footnote{或用该论文的术语，$F$-\Term{torsor}{挠子}。}当 $F$ 在 $X$ 上\Term{affine}{仿射}时，SGA 1 第 VIII 讲 2.1 还表明（仍用标准论证，参见 J. Giraud 的学位论文），$\mathrm H^1(X,F)$ 也是 $X$ 上\Term{scheme}{概形} $P$ 的类所成的群，其中 $F$ 在 $P$ 上\Term{right-action}{右作用}，且 $P$ 是“\Term{etale-topology}{平展拓扑}意义下 $X$ 上的\Term{principal-homogeneous-bundle}{主齐性纤维丛}”，即在\Term{etale-topology}{平展拓扑}意义下\Term{locally-trivial}{局部平凡}。\footnote{或用现在通行的术语，\Term{representable}{可表}的 $F$-\Term{torsor}{挠子}。}

\textbf{注记 2.1。} 须注意，若对每个 $X$ 上的\Term{scheme}{概形} $Z$，以 $\varphi_X(Z)$ 表示 $X_{\text{ét}}$ 上对应的\Term{sheaf}{层}，而 $f:Y\to X$ 是\Term{scheme}{概形}\Term{morphism}{态射}，则有一个显然的、关于 $Z$ \Term{functorial}{函子性}的\Term{homomorphism}{同态}
\[
 f^*(\varphi_X(Z))\longrightarrow\varphi_Y(Z_Y),\qquad Z_Y=Z\times_XY,\tag{*}
\]
也就是一个显然的\Term{homomorphism}{同态}
\[
 \varphi_X(Z)\longrightarrow f_*(\varphi_Y(Z_Y)),
\]
即关于 $X'\in\operatorname{Ob}X_{\text{ét}}$ \Term{functorial}{函子性}的\Term{homomorphism}{同态}
\[
 \operatorname{Hom}_X(X',Z)\longrightarrow\operatorname{Hom}_Y(X'_Y,Z_Y).
\]
但须注意，一般（$*$）\emph{不是}\Term{isomorphism}{同构}，即相对\Term{scheme}{概形}的对应\Term{etale}{平展}\Term{sheaf}{层}的形成\emph{不与\Term{inverse-image-functor}{逆像函子}可交换}。不过，在 $Z$ 于 $X$ 上\Term{etale}{平展}这一特殊情形中，（$*$）为\Term{isomorphism}{同构}。同样，当 $X\ne\varnothing$ 时，\Term{functor}{函子}
\[
 \varphi_X:\mathrm{Sch}_{/X}\longrightarrow\widetilde X_{\text{ét}}
\]
不\Term{faithful}{忠实}；\footnote{对此参见原书第 234 页页下注。}但其在 $X_{\text{ét}}$ 上的\Term{restriction}{限制}，即\Term{canonical-functor}{典范函子} $X_{\text{ét}}\to\widetilde X_{\text{ét}}$，是\Term{fully-faithful}{全忠实的}，因为由 a），$X_{\text{ét}}$ 的\Term{topology}{拓扑}比其\Term{canonical-topology}{典范拓扑}粗。

b）容易核实，（$*$）还与由 $I\in\mathcal U$ 编号的\Term{direct-sum}{直和}可交换。特别地，对每个\Term{set}{集合} $I$，以 $I_X$ 表示由 $I$ 定义的常值 $X$-\Term{scheme}{概形}（SGA 3 第 I 讲 1.8），即 $I$ 个 $X$ 的\Term{direct-sum}{直和}，则对应\Term{sheaf}{层}就是\Term{constant-sheaf}{常值层} $I_{X_{\text{ét}}}$，即 $I$ 个 $X_{\text{ét}}$ 上\Term{final-sheaf}{终层}的\Term{direct-sum}{直和}。因\Term{functor}{函子} $I\mapsto I_X$ 也与\Term{finite}{有限}\Term{projective-limit}{投射极限}可交换，可见它把群变为群、\Term{abelian-group}{交换群}变为\Term{abelian-group}{交换群}，等等。若 $G$ 为普通\Term{abelian-group}{交换群}，则把 $\mathrm H^i(X,G_X)$ 简写为 $\mathrm H^i(X,G)$。例如，假定 $G$ \Term{finite}{有限}，则 $G_X$ 在 $X$ 上\Term{finite}{有限}，因而\Term{affine}{仿射}，所以利用 a）末尾的注记，可把 $\mathrm H^1(X,G)$ 解释为群为 $G$ 的
\SourcePage{242}
\Term{principal-galois-covering}{主 Galois 覆盖}的类集（SGA 1 第 V 讲 2.7）。当 $X$ \Term{connected}{连通}并配备一个\Term{geometric-point}{几何点} $a$ 时，用“\Term{fundamental-group}{基本群}”$\pi_1(X,a)$（SGA 1 第 V 讲 7）可得\Term{canonical-isomorphism}{典范同构}
\[
 \mathrm H^1(X,G)\simeq\operatorname{Hom}(\pi_1(X,a),G),\tag{**}
\]
这里假定 $G$ 交换。

可以说，从 Zariski \Term{cohomology}{上同调}转到\Term{etale-topology}{平展拓扑}时，“所做恰到好处”，使\Term{finite}{有限}常值系数群 $G$ 的“正确的”$\mathrm H^1$（即（$*$）右端所出现的群）得以取得。一个值得重视的事实是，这也足以对每个\Term{torsion-group}{挠系数群}求得“正确的”$\mathrm H^i(X,G)$，至少当 $G$ 与 $X$ 的各\Term{residue-characteristic}{剩余特征}互素时如此；本研讨班后面将证明这一事实。

c）设 $F$ 是 $X$ 上 Zariski \Term{topology}{拓扑}意义下的 $\mathscr O_X$-模。用 SGA 3 第 I 讲 4.6 的记号，定义\Term{functor}{函子}
\[
 W(F):(\mathrm{Sch})_{/X}^\circ\longrightarrow(\mathrm{Ens}),\qquad
 W(F)(X')=\Gamma(X',F\otimes_{\mathscr O_X}\mathscr O_{X'}).
\]
由 SGA 1 第 VIII 讲 1.6 与 SGA 3 第 IV 讲 6.3.1，这个\Term{functor}{函子}仍为\Term{etale-topology}{平展拓扑}（甚至 \Term{fpqc-topology}{fpqc 拓扑}）的\Term{sheaf}{层}。更有甚者，$W(F)$ 在 $X_{\text{ét}}$ 上的\Term{restriction}{限制}仍为\Term{sheaf}{层}，仍记作 $W(F)$。按定义，有 $\mathrm H^0(X,W(F))=\Gamma(F)=\mathrm H^0(X,F)$。当 $F$ \Term{quasi-coherent}{拟凝聚}时，还有更强的结果，参见 4.3。

\SGAsection{3}{平展拓扑斯的生成元；层的归纳极限的上同调}
\begin{SGAresult}{命题}{3.1}
设 $X$ 为\Term{scheme}{概形}，$C$ 为 $X_{\text{ét}}$ 的一个\Term{full-subcategory}{全子范畴}，使对每个在 $X$ 上\Term{etale}{平展}且\Term{affine}{仿射}的 $X'$，$C$ 都含有一个与 $X'$ \Term{isomorphism}{同构}的\Term{object}{对象}。则 $C$ 是\Term{site}{景} $X_{\text{ét}}$ 的“\Term{topological-generators}{拓扑生成元族}”，因而是\Term{topos}{拓扑斯} $\widetilde X_{\text{ét}}$ 的\Term{generators}{生成元族}，且\Term{restriction}{限制}\Term{functor}{函子} $\widetilde X_{\text{ét}}\to\widetilde C$ 是\Term{equivalence-of-categories}{范畴等价}，其中 $C$ 配备由 $X_{\text{ét}}$ 诱导的\Term{topology}{拓扑}。
\end{SGAresult}
\begin{proof}
借助“\Term{comparison-lemma}{比较引理}”，结论是平凡的。
\end{proof}

\begin{SGAresult}{推论}{3.2}
假定 $X$ \Term{quasi-separated}{拟分离}。把 $X_{\text{ét}}$ 中由在 $X$ 上\Term{finite-presentation}{有限表现}的\Term{etale}{平展}\Term{scheme}{概形}组成的\Term{full-subcategory}{全子范畴} $C$，配备由 $X_{\text{ét}}$ 诱导的\Term{topology}{拓扑}，称为 $X$ 的\Term{restricted-etale-site}{限制平展景}。则：

（i）$C$ 对\Term{fiber-product}{纤维积}\Term{stable}{稳定}；若 $X$ \Term{quasi-compact}{拟紧}，则 $C$ 是\Term{finite-type-site}{有限型景}。

（ii）\Term{restriction}{限制}\Term{functor}{函子} $\widetilde X_{\text{ét}}\to\widetilde C$ 是\Term{equivalence-of-categories}{范畴等价}。
\end{SGAresult}
\begin{proof}
断言（i）平凡；（ii）由 3.1 得出，因为 $X$ \Term{quasi-separated}{拟分离}保证 $C$ 满足 3.1 的条件：若 $X'$ \Term{quasi-compact}{拟紧}，
\SourcePage{243}
例如\Term{affine}{仿射}，则它在 $X$ 上\Term{quasi-compact}{拟紧}（EGA IV 1.2.4，取 $Z=\operatorname{Spec}\mathbf Z$），故若 $X'$ 在 $X$ 上\Term{locally-finite-type}{局部有限型}，便在 $X$ 上\Term{finite-presentation}{有限表现}。
\end{proof}

\begin{SGAresult}{命题}{3.3}
假定 $X$ \Term{quasi-compact}{拟紧}且\Term{quasi-separated}{拟分离}。则 $X_{\text{ét}}$ 上\Term{abelian-sheaf}{阿贝尔层}\Term{category}{范畴}上的\Term{functor}{函子} $\mathrm H^q(X_{\text{ét}},F)$ 与\Term{inductive-limit}{归纳极限}可交换。
\end{SGAresult}
\begin{proof}
由 3.2（ii），可用 $C$ 代替 $X_{\text{ét}}$；因 $X$ \Term{quasi-compact}{拟紧}，有 $X\in\operatorname{Ob}C$，结论由 3.2（i）与第 VI 讲 6.1.2（3）得出。
\end{proof}

\SGAsection{4}{与其他拓扑的比较}
\textbf{4.0.} 首先注意，第 2 节所考虑的 $X_{\text{ét}}$ 上\Term{sheaf}{层}的例子，自然地都是定义在配备\Term{etale-topology}{平展拓扑}（甚至 \Term{fpqc-topology}{fpqc 拓扑}）的 $(\mathrm{Sch})_{/X}$ 上的\Term{sheaf}{层}的\Term{restriction}{限制}。一般地，设
\[
 u^*:X_{\text{ét}}\longrightarrow(\mathrm{Sch})_{/X}
\]
为\Term{inclusion-functor}{包含函子}；它\Term{continuous-functor}{连续}（第 III 讲）且与\Term{finite}{有限}\Term{projective-limit}{投射极限}可交换，因而定义\Term{site-morphism}{景态射}
\[
 u:(\mathrm{Sch})_{/X}\longrightarrow X_{\text{ét}},
\]
从而在相应的\Term{sheaf}{层}\Term{category}{范畴}上给出\Term{functor}{函子}
\[
 u_*:\widetilde{(\mathrm{Sch})_{/X}}\longrightarrow\widetilde X_{\text{ét}},
 \qquad u_*(F)\longrightarrow F\circ u,
\]
以及后者的一个左\Term{adjoint-functor}{伴随函子}\footnote{严格说来，$(\mathrm{Sch})_{/X}$ 是 $\mathcal U$-\Term{category}{范畴}而不是 $\mathcal U$-\Term{site}{景}，所以不能援引第 III 讲来保证 $u^*$ 存在。不过，容易通过公式
\[
 u^*(F)(X')=(p_{X'})_{\text{ét}}^*(F)(X')\qquad
 (X'\in\operatorname{Ob}\mathrm{Sch}_{/X})
\]
构造 $u^*$，其中 $p_{X'}$ 表示\Term{structural-morphism}{结构态射} $X'\to X$。另一方面，要使 4.1（ii）有意义，并使所述 4.1 的证明成立，应引入满足 $\mathcal U\in\mathcal V$ 的\Term{universe}{宇宙} $\mathcal V$，把公式 4.1（i）中出现的\Term{sheaf}{层}看作取值于 $\mathcal V$-$(\mathrm{Ens})$ 的\Term{sheaf}{层}。若只想使用 $\mathcal U$-\Term{site}{景}，也可把 $(\mathrm{Sch})_{/X}$ 换成它的一个 $\mathcal U$-小\Term{full-subcategory}{全子范畴}，使之对\Term{finite}{有限}\Term{projective-limit}{投射极限}\Term{stable}{稳定}且包含 3.1 的 $C$。}
\[
 u^*:\widetilde X_{\text{ét}}\longrightarrow\widetilde{\mathrm{Sch}_{/X}}.
\]
于是：

\begin{SGAresult}{命题}{4.1}
（i）\Term{functor}{函子} $u^*$ \Term{fully-faithful}{全忠实}，因而对 $X_{\text{ét}}$ 上每个\Term{sheaf}{层} $G$，\Term{canonical-homomorphism}{典范同态} $G\to u_*u^*(G)$ 是\Term{isomorphism}{同构}。

（ii）设 $G$（相应地，$F$）为 $X_{\text{ét}}$（相应地，$(\mathrm{Sch})_{/X}$）上的\Term{abelian-sheaf}{阿贝尔层}。则下列\Term{canonical-homomorphism}{典范同态}（例如可定义为 Leray \Term{spectral-sequence}{谱序列}的\Term{edge-homomorphism}{边缘同态}）是\Term{isomorphism}{同构}：
\SourcePage{244}
\[
 \mathrm H^*(X_{\text{ét}},F\circ u)\xrightarrow\sim\mathrm H^*((\mathrm{Sch})_{/X},F),
\]
\[
 \mathrm H^*(X_{\text{ét}},F\circ u)\xrightarrow\sim\mathrm H^*((\mathrm{Sch})_{/X},F).
\]
\end{SGAresult}
\begin{proof}
事实上，\Term{functor}{函子} $u$ 是一个\Term{comorphism}{余态射}，因而可应用第 III 讲。
\end{proof}
所以，（i）表明，为研究\Term{sheaf}{层}的\Term{cohomology}{上同调}，在“小”\Term{etale-site}{平展景} $X_{\text{ét}}$ 上工作与在“大”\Term{etale-site}{平展景} $(\mathrm{Sch})_{/X}$ 上工作，本质上是\Term{equivalence}{等价}的。

\textbf{4.2.} 另一方面，有必要在 $(\mathrm{Sch})$（从而在各 $(\mathrm{Sch})_{/X}$）中引入\Term{etale-topology}{平展拓扑}以外的其他\Term{topology}{拓扑}；其中最有用的定义在 SGA 3 第 IV 讲 6.3 中。它们中最粗的是 Zariski \Term{topology}{拓扑}（Zar），由\Term{covering-family}{覆盖族}为\Term{open-immersion}{开浸入}的\Term{surjective-family}{满射族}的\Term{pretopology}{预拓扑}定义；它比\Term{etale-topology}{平展拓扑}粗。这些\Term{topology}{拓扑}中最细的是“\Term{faithfully-flat}{忠实平坦}\Term{quasi-compact}{拟紧}”\Term{topology}{拓扑}，简称 fpqc；它是使 Zariski 意义下的\Term{covering-family}{覆盖族}及\Term{faithfully-flat}{忠实平坦}\Term{quasi-compact}{拟紧}\Term{morphism}{态射}均为\Term{covering}{覆盖}的最\Term{coarse-topology}{粗拓扑}。\Term{fpqc-topology}{fpqc 拓扑}比\Term{etale-topology}{平展拓扑}细。如前所述，第 2 节所考虑的 $(\mathrm{Sch})_{/X}$ 上\Term{sheaf}{层}的各种例子，实际上已经是 \Term{fpqc-topology}{fpqc 拓扑}的\Term{sheaf}{层}。

\textbf{4.2.1.} 不过须注意，对于 $(\mathrm{Sch})_{/X}$ 上 \Term{fpqc-topology}{fpqc 拓扑}（或 SGA 3 中考虑的 fppf 等\Term{topology}{拓扑}）的\Term{abelian-sheaf}{阿贝尔层} $F$，$F$ 的 \Term{fpqc-topology}{fpqc 拓扑}（相应地，……）\Term{cohomology-group}{上同调群}并不总\Term{isomorphism}{同构}于\Term{etale-topology}{平展拓扑}的\Term{cohomology-group}{上同调群}；即使 $X$ 是一个域 $k$ 的谱、$F$ 由 $k$ 上的一个\Term{algebraic-group}{代数群}表示，甚至对 $\mathrm H^1$ 也如此。一般可以证明，对于系数群为 $X$ 上\Term{etale}{平展}、或更一般\Term{smooth}{光滑}的\Term{group-scheme}{群概形}的情形，\Term{etale-topology}{平展拓扑}给出“正确的”\Term{cohomology-group}{上同调群}。\footnote{参见 A. Grothendieck，《Brauer 群》III，定理 11.7，载《\Term{scheme}{概形}\Term{topology}{拓扑}十讲》（Dix exposés sur la topologie des schémas），North Holland Pub. Cie。}但对于诸如 $S$ 上的\Term{radicial}{根性}群这样的\Term{group-scheme}{群概形}，便不再如此；此时应将仍然太粗的\Term{etale-topology}{平展拓扑}换成 fpqc 或 fppf \Term{topology}{拓扑}。

\textbf{4.2.2.} 作为不同\Term{topology}{拓扑}的\Term{cohomology}{上同调}之间关系的例子，这里指出 Zariski \Term{topology}{拓扑}与\Term{etale-topology}{平展拓扑}的情形。以 $X_{\mathrm{Zar}}$ 表示 $X$ 的 Zariski \Term{open-subset}{开集}\Term{site}{景}，于是有\Term{canonical}{典范}\Term{inclusion-functor}{包含函子}
\[
 u^*:X_{\mathrm{Zar}}\longrightarrow X_{\text{ét}},
\]
它定义\Term{site-morphism}{景态射}
\[
 u:X_{\text{ét}}\longrightarrow X_{\mathrm{Zar}},
\]
\SourcePage{245}
从而有相应的\Term{direct-image-functor}{正像函子}
\[
 f_*:\widetilde X_{\text{ét}}\longrightarrow\widetilde X_{\mathrm{Zar}},
 \qquad f_*(F)=F\circ u,
\]
及\Term{inverse-image-functor}{逆像函子}
\[
 f^*:\widetilde X_{\mathrm{Zar}}\longrightarrow\widetilde X_{\text{ét}},
\]
二者互为\Term{adjunction}{伴随}；从几何上，应把偶对 $(f_*,f^*)$ 理解为一个\Term{topos}{拓扑斯}\Term{morphism}{态射}
\[
 f:\widetilde X_{\text{ét}}\longrightarrow\widetilde X_{\mathrm{Zar}}.
\]
由此得到\Term{cohomological-functor}{上同调函子}的\Term{homomorphism}{同态}
\[
 \mathrm H^*(X_{\mathrm{Zar}},f_*(F))\longrightarrow\mathrm H^*(X_{\text{ét}},F)
\]
以及 Leray \Term{spectral-sequence}{谱序列}
\[
 \mathrm H^*(X_{\text{ét}},F)\Longleftarrow
 \mathrm H^p(X_{\mathrm{Zar}},\mathrm R^qf_*(F)),
\]
其中 $F$ 是 $X_{\text{ét}}$ 上的\Term{abelian-sheaf}{阿贝尔层}。这个\Term{spectral-sequence}{谱序列}概括了\Term{etale-cohomology}{平展上同调}与 Zariski \Term{cohomology}{上同调}之间的一般关系。对于\Term{constant-sheaf}{常值系数层}等系数\Term{sheaf}{层} $F$，这个\Term{spectral-sequence}{谱序列}一般当然远非平凡，即一般有 $\mathrm R^qf_*(F)\ne0$（尤其可取 $X$ 为一个域的谱）。然而：

\begin{SGAresult}{命题}{4.3}
设 $F$ 为\Term{prescheme}{预概形} $X$ 上可变的\Term{sheaf-of-modules}{模层}（Zariski \Term{topology}{拓扑}意义下的\Term{sheaf}{层}），从而得到 $X_{\text{ét}}$ 上的\Term{sheaf}{层} $F_{\text{ét}}$（参见第 2 节 c））及\Term{cohomological-functor}{上同调函子}的\Term{homomorphism}{同态}
\[
 \mathrm H^*(X,F)\longrightarrow\mathrm H^*(X_{\text{ét}},F_{\text{ét}}).
\]
当 $F$ \Term{quasi-coherent}{拟凝聚}时，上述\Term{homomorphism}{同态}是\Term{isomorphism}{同构}。
\end{SGAresult}
\begin{proof}
事实上，左端就是 $\mathrm H^*(X_{\mathrm{Zar}},f_*(F_{\text{ét}}))$；用前面的记号，只须证明
\[
 \mathrm R^qf_*(F)=0\qquad(q>0).
\]
由 $\mathrm R^qf_*$ 的计算方法，以及\Term{affine}{仿射}\Term{open-subset}{开集}构成 $X$ \Term{topology}{拓扑}的一组基这一事实，这立即归结到下面的推论。
\end{proof}

\begin{SGAresult}{推论}{4.4}
若 $X$ \Term{affine}{仿射}，则
\[
 \mathrm H^q(X_{\text{ét}},F_{\text{ét}})=0\qquad(q>0).
\]
\end{SGAresult}
\begin{proof}
为说明这一点，设 $C=X_{\text{ét}}$，$C'$ 为由\Term{affine}{仿射}的 $X'$ 组成的\Term{full-subcategory}{全子范畴}。由 3.1，
\[
 \mathrm H^q(C,F_{\text{ét}})\simeq\mathrm H^q(C',F|_{C'}).
\]
所以只须证明 $F|_{C'}$ 是 $C'$-\Term{acyclic}{无上同调}的（第 V 讲 4.1），或\Term{equivalence}{等价}地满足第 V 讲 4.3 的条件，即对每个 $X'\in\operatorname{Ob}C'$ 和 $C'$ 中的每个\Term{covering-family}{覆盖族} $\mathscr R=(X'_i\to X')_i$，有 $\mathrm H^q(\mathscr R,F)=0$，$q>0$。可以假定 $(X'_i)_i$ \Term{finite}{有限}，
\SourcePage{246}
再将各 $X'_i$ 换成它们的和，可设\Term{covering-family}{覆盖族}只含一个\Term{covering}{覆盖}\Term{morphism}{态射} $X'_i\to X'$，即（\Term{etale}{平展}且）\Term{surjectivity}{满射}的\Term{morphism}{态射}。于是归结为证明，若 $f:X\to Y$ 是\Term{affine}{仿射}\Term{scheme}{概形}的\Term{etale}{平展}\Term{surjectivity}{满射}\Term{morphism}{态射}，而 $F$ 是 $X$ 上的\Term{quasi-coherent}{拟凝聚}模，则
\[
 \mathrm H^q(X/Y,F)=0\qquad(q>0).\tag{*}
\]
这在 TDTE I，B，1.1 中已经证明（载于 FGA，参见第 IX 讲文献 [4]）。
\end{proof}

\textbf{注记 4.5。} 事实上，由该处结果，要得到（$*$），只须 $X\to Y$ \Term{surjectivity}{满射}且\Term{flat}{平坦}（不必\Term{etale}{平展}）。因此，上述证明还表明，把 3.3 中的\Term{etale-site}{平展景} $X_{\text{ét}}$ 换成配备 SGA 3 第 IV 讲 6.3 中所考虑的任何比 Zariski \Term{topology}{拓扑}更细的\Term{topology}{拓扑}（例如 \Term{fpqc-topology}{fpqc 拓扑}）的 $(\mathrm{Sch})_{/X'}$，仍有类似的\Term{isomorphism}{同构}。

\SGAsection{5}{概形投射极限的上同调}
\textbf{5.1.} 设 $I$ 是一个递增的\Term{filtered-preordered-set}{滤过预序集}，$\mathfrak X=(X)_{i\in I}$ 为\Term{scheme}{概形}的一个\Term{projective-system}{投射系统}，其\Term{transition-morphism}{转移态射} $u_i:X_i\to X_j$（$i\geq j$）均\Term{affine}{仿射}。回顾（EGA IV 8……），在这些条件下，\Term{projective-limit}{投射极限}
\[
 X=\varprojlim X_i
\]
存在于\Term{scheme}{概形}\Term{category}{范畴}中（甚至存在于\Term{locally-ringed-space}{局部赋环空间}\Term{category}{范畴}中），且可如下构造：选择 $i_0\in I$，于是对 $i\geq i_0$ 有一个 $X_{i_0}$-\Term{isomorphism}{同构}（EGA II 1.2 和 1.3）
\[
 X_i=\operatorname{Spec}\mathscr A_i,
\]
其中 $\mathscr A_i$ 是 $X_i$ 上的\Term{quasi-coherent}{拟凝聚}\Term{algebra}{代数}；各 $\mathscr A_i$（$i\geq i_0$）遂构成这种\Term{algebra}{代数}的一个\Term{inductive-system}{归纳系统}。置
\[
 \mathscr A=\varinjlim_i\mathscr A_i,
\]
则可取 $X=\operatorname{Spec}\mathscr A$。此外，若以 $\operatorname{esp}S$ 表示\Term{scheme}{概形} $S$ 的底\Term{topological-space}{拓扑空间}，可以证明（EGA IV 8……），\Term{canonical}{典范}映射
\[
 \operatorname{esp}X\longrightarrow\varprojlim\operatorname{esp}(X_i)
\]
是\Term{homeomorphism}{同胚}，而\Term{sheaf-of-rings}{环层}的\Term{canonical-homomorphism}{典范同态}
\[
 \varinjlim_i\gamma_i^{-1}(\mathscr O_{X_i})\longrightarrow\mathscr O_X
\]
是\Term{bijection}{双射}，其中 $\gamma_i:X\to X_i$ 为\Term{canonical}{典范}\Term{continuous-map}{连续映射}。

\SourcePage{247}
\textbf{5.2.} 形象地说，EGA IV 8、9 的内容可以概括为：当 $X_{i_0}$ \Term{quasi-compact}{拟紧}且\Term{quasi-separated}{拟分离}时，“$X$ 上每个在 $X$ 上\Term{finite-presentation}{有限表现}的\Term{scheme}{概形}性质的资料”，都\Term{equivalence}{等价}于“$i$ \Term{large-relative-to-cardinal}{足够大}时”某个 $X_i$ 上的同类资料。例如，可证明（EGA IV 8……）：

a）若 $i_1\in I$，而 $X'_{i_1},X''_{i_1}$ 是 $X_{i_1}$ 上两个\Term{finite-presentation}{有限表现}\Term{scheme}{概形}，对每个 $i\geq i_1$ 置
\[
 X'_i=X'_{i_1}\times_{X_{i_1}}X_i,\qquad
 X'=X'_{i_1}\times_{X_{i_1}}X,
\]
并类似地定义 $X''_i,X''$，则\Term{canonical}{典范}映射
\[
 \varinjlim_{i\geq i_1}\operatorname{Hom}_{X_i}(X'_i,X''_i)
 \longrightarrow\operatorname{Hom}_X(X',X'')
\]
是\Term{bijection}{双射}。

b）对每个 $X$ 上的\Term{finite-presentation}{有限表现}\Term{scheme}{概形} $X'$，存在指标 $i_1\in I$、$X_{i_1}$ 上的\Term{finite-presentation}{有限表现}\Term{scheme}{概形} $X'_{i_1}$，以及一个 $X$-\Term{isomorphism}{同构}
\[
 X'\simeq X'_{i_1}\times_{X_{i_1}}X.
\]

\textbf{5.3.} 可以用第 VI 讲引入的\Term{fibered-category}{纤维范畴}的\Term{inductive-limit}{归纳极限}语言表述上述结果，以及 EGA IV 8 中性质相似的结果。为此，考虑由\Term{scheme}{概形}的\Term{finite-presentation}{有限表现}\Term{morphism}{态射} $f:T\to S$ 组成的\Term{category}{范畴} $\mathscr F$，及“取靶\Term{functor}{函子}”
\[
 \pi:\mathscr F\longrightarrow(\mathrm{Sch}),
\]
它显然是一个\Term{fibered-functor}{纤维化函子}，而且各\Term{fiber-category}{纤维范畴}都\Term{equivalence}{等价}于 $\mathcal U$-\Term{small-category}{小范畴}。考虑\Term{fibered-category}{纤维范畴} $\mathscr F/(\mathrm{Sch})$ 沿\Term{functor}{函子}
\[
 i\longmapsto X_i:\operatorname{cat}(I)\longrightarrow(\mathrm{Sch})
\]
的\Term{inverse-image}{逆像}，其中 $\operatorname{cat}(I)$ 是与 $I$ 关联的\Term{category}{范畴}，$\operatorname{Hom}(i,j)\ne\varnothing$ 当且仅当 $i\geq j$。这给出一个\Term{fibered-category}{纤维范畴}
\[
 \pi_{\mathfrak X}:\mathscr F_{\mathfrak X}\longrightarrow\operatorname{cat}(I),
\]
其在每个 $i\in I$ 处的\Term{fiber}{纤维}\Term{canonical-isomorphism}{典范同构}于 $(\mathrm{Sch})_{/X_i}$ 中由 $X_i$ 上\Term{finite-presentation}{有限表现}\Term{scheme}{概形}组成的\Term{full-subcategory}{全子范畴} $\mathscr F_{X_i}$；对应于 $j\geq i$ 的\Term{base-change}{基变换}\Term{morphism}{态射} $\mathscr F_i\to\mathscr F_j$ 正是 $X'_i\mapsto X'_j=X'_i\times_{X_i}X_j$。于是，把每个 $X_i$ 上\Term{finite-presentation}{有限表现}的 $X'_i$ 送到其在 $X$ 上的\Term{inverse-image}{逆像}
\[
 \varphi(X'_i)=X'_i\times_{X_i}X,
\]
便得到一个自然\Term{functor}{函子}
\[
 \varphi:\mathscr F_{\mathfrak X}\longrightarrow\mathscr F_{X'},
\]
\SourcePage{248}
它显然把 $\mathscr F_{\mathfrak X}$ 的\Term{cartesian-morphism}{笛卡尔态射}变成\Term{isomorphism}{同构}，因此按定义，它唯一地通过一个\Term{canonical-functor}{典范函子}\Term{resolution}{分解}：
\[
 \psi:\varinjlim_{\mathscr F_{\mathfrak X}/\operatorname{Cat}I}\mathscr F_{\mathfrak X}
 \longrightarrow\mathscr F_X.\tag{*}
\]
结合左端的描述，上面回顾的 EGA IV 8 的结果即可表述为：（$*$）是一个\Term{equivalence-of-categories}{范畴等价}。

\textbf{5.4.} 可以用多种方式细化上述 5.2 a）的断言：引入一个对\Term{base-change}{基变换}\Term{stable}{稳定}的\Term{scheme}{概形}\Term{morphism}{态射}\Term{set}{集合} $(\mathcal M)$，并断言，对给定的 $X_{i_1}$-\Term{morphism}{态射} $u_{i_1}:X'_{i_1}\to X''_{i_1}$，相应\Term{morphism}{态射} $u':X'\to X''$ 属于 $\mathcal E(\mathcal M)$，当且仅当存在 $i\geq i_1$，使 $X_i$-\Term{morphism}{态射} $u_i:X'_i\to X''_i$ 属于 $\mathcal E(\mathcal M)$。例如，当 $(\mathcal M)$ 是以下任一个\Term{arrow}{箭头}\Term{set}{集合}时，所得断言成立：\Term{proper}{固有}\Term{morphism}{态射}（相应地，\Term{projective-morphism}{射影}、\Term{quasi-projective}{拟射影}、\Term{affine}{仿射}、\Term{affine}{仿射}、\Term{finite-morphism}{有限}、\Term{quasi-finite}{拟有限}、\Term{radicial}{根性}、\Term{surjectivity}{满射}、\Term{flat}{平坦}、\Term{smooth}{光滑}、\Term{etale}{平展}、\Term{unramified}{非分歧}\Term{morphism}{态射}）。读者可在 EGA IV 第 8 节（前九种情形）、第 11 节（\Term{flat}{平坦}情形）和第 17 节（最后三种情形）中找到相应的断言。

\textbf{5.5.} 现在把 $\mathscr F\to(\mathrm{Sch})$ 换成由\Term{finite-presentation}{有限表现}\Term{etale-morphism}{平展态射}组成的\Term{fibered-subcategory}{纤维子范畴} $\mathscr G$；对每个\Term{scheme}{概形} $S$，其\Term{fiber-category}{纤维范畴} $\mathscr G_S$ 就是 $S$ 上\Term{finite-presentation}{有限表现}\Term{etale}{平展}\Term{scheme}{概形}的\Term{category}{范畴}。类似地，在 $\operatorname{cat}(I)$ 上构造\Term{fibered-category}{纤维范畴} $\mathscr G_{\mathfrak X}$，它在每个 $i\in I$ 处的\Term{fiber-category}{纤维范畴}为 $\mathscr G_{X_i}$。由 3.2（ii），\Term{etale-topos}{平展拓扑斯} $\widetilde X_{i,\text{ét}}$、$\widetilde X_{\text{ét}}$ 还分别\Term{canonical}{典范}\Term{equivalence}{等价}于 $\widetilde{\mathscr G}_{X_i}$、$\widetilde{\mathscr G}_X$，其中每个 $\mathscr G_S$ 都配备由 $S_{\text{ét}}$ 诱导的\Term{topology}{拓扑}。此外，由 3.2（i），每个 $\mathscr G_{X_i}$ 都是一个\Term{finite-type-site}{有限型景}，而且\Term{equivalence}{等价}于一个 $\mathcal U$-小\Term{site}{景}。考虑各 $\mathscr G_{X_i}$ 上的\Term{topology}{拓扑}后，$\mathscr G$ 成为 $\operatorname{cat}(I)$ 上的一个\Term{fibered-site}{纤维景}。于是可陈述：

\begin{SGAresult}{引理}{5.6}
\Term{canonical-functor}{典范函子}
\[
 \varinjlim_{\mathscr G_{\mathfrak X}/\operatorname{cat}(I)}\mathscr G_{\mathfrak X}
 \longrightarrow\mathscr G_X
\]
定义一个保持\Term{topology}{拓扑}的\Term{equivalence}{等价}，把各 $X_i$ 的\Term{restricted-etale-site}{限制平展景}（3.2）的\Term{inductive-limit}{归纳极限}\Term{site}{景}，等同于 $X=\varprojlim_iX_i$ 的\Term{restricted-etale-site}{限制平展景}。
\end{SGAresult}
\begin{proof}
所得\Term{functor}{函子}是\Term{equivalence-of-categories}{范畴等价}，这是 5.4 所回顾的断言之一，即 $(\mathcal M)$ 为 $(\mathrm{Sch})$ 中\Term{etale-morphism}{平展态射}\Term{set}{集合}的情形。关于\Term{topology}{拓扑}的断言，同样可取 $(\mathcal M)$ 为 $(\mathrm{Sch})$ 中\Term{surjectivity}{满射}\Term{morphism}{态射}的\Term{set}{集合}而得到。
\end{proof}
上述结果使我们可将\SGAsourcegap{原文在“les résultats de”之后缺少引用}的结果应用于目前的情形。特别得到：

\SourcePage{249}
\begin{SGAresult}{定理}{5.7}
设 $\mathfrak X=(X_i)_{i\in I}$ 为\Term{scheme}{概形}的\Term{projective-system}{投射系统}，$I$ 为递增\Term{filtered-preordered-set}{滤过预序集}。假定\Term{transition-morphism}{转移态射} $u_{ij}:X_j\to X_i$ \Term{affine}{仿射}（所以\Term{scheme}{概形} $X=\varprojlim X_i$ 有定义（5.1）），且各 $X_i$ \Term{quasi-compact}{拟紧}并\Term{quasi-separated}{拟分离}。考虑 5.5 中具体给出的 $\operatorname{cat}(I)$ 上的\Term{fibered-site}{纤维景} $\mathscr G$，它在 $i\in I$ 处的\Term{fiber}{纤维} $\mathscr G_i$ \Term{canonical-isomorphism}{典范同构}于 $X_i$ 的“\Term{restricted-etale-site}{限制平展景}”（由 $X_i$ 上\Term{finite-presentation}{有限表现}\Term{etale}{平展}\Term{scheme}{概形}组成）。设 $F$ 是 $\mathscr G$ 上的\Term{abelian-sheaf}{阿贝尔层}，即（3）一个\Term{functor}{函子} $\mathscr G^\circ\to(\mathrm{Ab})$，它在每个 $\mathscr G_i$ 上的\Term{restriction}{限制} $F_i$ 都是\Term{sheaf}{层}，亦即 $X_i$ 上的\Term{sheaf}{层}。设 $F_\infty=\varinjlim_i u_i^*(F_i)$ 为 $X$ 上的诱导\Term{sheaf}{层}，其中 $u_i:X\to X_i$ 为\Term{canonical}{典范}\Term{morphism}{态射}。则\Term{canonical-homomorphism}{典范同态}
\[
 \varinjlim_i\mathrm H^n(X_i,F_i)\longrightarrow\mathrm H^n(X,F_\infty)
\]
都是\Term{isomorphism}{同构}。
\end{SGAresult}
\begin{proof}
结合 5.6，这是第 VI 讲 8.7.3 与第 VI 讲 8.5.2 的共同结果。
\end{proof}
我们主要将在以下特殊情形中使用 5.7，将其明确写为：

\begin{SGAresult}{推论}{5.8}
对 $(X_i)_{i\in I}$、$X$、$u_{ij}$、$u_i$ 采用前面的假设与记号，设 $i_0\in I$，$F_{i_0}$ 是 $X_{i_0}$ 上的\Term{abelian-sheaf}{阿贝尔层}。对每个 $i\geq i_0$，置 $F_i=u_{i_0i}^*(F_{i_0})$，再置 $F_\infty=u_i^*(F_{i_0})$。在这些条件下，\Term{canonical-homomorphism}{典范同态}
\[
 \varinjlim_i\mathrm H^n(X_i,F_i)\longrightarrow\mathrm H^n(X,F_\infty)
\]
都是\Term{isomorphism}{同构}。
\end{SGAresult}
以下是一个有时有用、属于 5.7 而不属于 5.8 的情形：

\begin{SGAresult}{推论}{5.9}
对 $(X_i)_{i\in I}$、$X$、$u_{ij}$、$u_i$ 采用 5.7 的假设与记号。给定 $i_0\in I$，设 $G_{i_0}$ 为 $X_{i_0}$ 上\Term{locally-finite-presentation}{局部有限表现}的交换\Term{group-scheme}{群概形}。对每个 $i\geq i_0$，置 $G_i=G_{i_0}\times_{X_{i_0}}X_i$，并类似地置 $G_\infty=G_{i_0}\times_{X_{i_0}}X$。则\Term{canonical-homomorphism}{典范同态}
\[
 \varinjlim_i\mathrm H^n(X_i,G_i)\longrightarrow\mathrm H^n(X,G_\infty)
\]
都是\Term{isomorphism}{同构}。
\end{SGAresult}
\begin{proof}
已知 $G_{i_0}$ 在 $(\mathrm{Sch})_{/X_{i_0}}$ 上定义一个\Term{sheaf}{层}（参见第 2 节 a））；可假定 $i_0$ 是 $I$ 的最小\Term{object}{对象}，于是有\Term{canonical-functor}{典范函子} $\mathscr G\to(\mathrm{Sch})_{/X_{i_0}}$，复合后便得到 $\mathscr G$ 上的一个\Term{contravariant-functor}{反变函子} $F$，其在 $\mathscr G_i$ 上的\Term{restriction}{限制}\Term{canonical-isomorphism}{典范同构}于 $G_i$ 在 $X_i$ 上定义的\Term{sheaf}{层}。特别地，$F$ 是 $\mathscr G$ 上的\Term{sheaf}{层}，可对它应用 5.7。假设 $G_{i_0}$ 在 $X_{i_0}$ 上\Term{locally-finite-presentation}{局部有限表现}，是为了保证 5.7 的\Term{sheaf}{层} $F_\infty$ 经自然\Term{homomorphism}{同态} $F_\infty\to\operatorname{faisc}(G_\infty)$ \Term{isomorphism}{同构}于 $G_\infty$ 定义的\Term{sheaf}{层} $\operatorname{faisc}(G_\infty)$；
\SourcePage{250}
事实上，考察所考虑的\Term{sheaf}{层}在 $G_\infty$ 的\Term{restricted-etale-site}{限制平展景}的\Term{object}{对象}上的值，并应用 5.6，这立即由 EGA IV 8.8.2（i）得出。5.9 的证明至此完成。
\end{proof}

\begin{SGAresult}{推论}{5.10}
对 $(X_i)_{i\in I}$、$X$、$u_{ij}$、$u_i$ 采用 5.7 的假设与记号。设 $F$ 为 $X$ 上的交换\Term{sheaf-of-groups}{群层}，则\Term{canonical-homomorphism}{典范同态}
\[
 \varinjlim_i\mathrm H^n(X_i,u_{i*}(F))\longrightarrow\mathrm H^n(X,F)
\]
都是\Term{isomorphism}{同构}。
\end{SGAresult}
借助第 VI 讲 8.7.9，断言 5.7 至 5.10 可以推广为\Term{sheaf-of-modules}{模层}的 $\operatorname{Ext}^i$ 的相应断言；这里不再把它们在目前特殊情形下的形式重复一遍。不过，为便于以后引用，下面用\Term{functor}{函子} $\mathrm R^nf_*$ 明确写出 5.8 和 5.9 的变体。

\begin{SGAresult}{推论}{5.11}
设 $I$ 为递增\Term{filtered-preordered-set}{滤过预序集}，$(X_i)_{i\in I}$ 与 $(Y_i)_{i\in I}$ 为两个\Term{transition-morphism}{转移态射}\Term{affine}{仿射}的\Term{scheme}{概形}\Term{projective-system}{投射系统}，$X=\varprojlim_iX_i$、$Y=\varprojlim_iY_i$；$(f_i)_{i\in I}$ 为\Term{quasi-compact}{拟紧}且\Term{quasi-separated}{拟分离}\Term{morphism}{态射} $f_i:X_i\to Y_i$ 的一个\Term{projective-system}{投射系统}。设 $i_0\in I$，$F_{i_0}$ 为 $X_{i_0}$ 上的\Term{sheaf}{层}。对每个 $i\geq i_0$，令 $F_i$ 为 $F_{i_0}$ 在 $X_i$ 上的\Term{inverse-image}{逆像}\Term{sheaf}{层}。同样，令 $f:X\to Y$ 由 $(f_i)$ \Term{derivation}{导出}，$F$ 为 $F_{i_0}$ 在 $X$ 上的\Term{inverse-image}{逆像}\Term{sheaf}{层}。在这些条件下，\Term{canonical-homomorphism}{典范同态}
\[
 \varinjlim_i u_i^*(\mathrm R^nf_{i*}(F_i))\longrightarrow\mathrm R^nf_*(F)
\]
是\Term{isomorphism}{同构}，其中 $u_i:X\to X_i$ 为\Term{canonical}{典范}\Term{morphism}{态射}。
\end{SGAresult}
\begin{proof}
可假定 $i_0$ 是 $I$ 的最小\Term{object}{对象}；由于问题在 $Y_{i_0}$ 上是\Term{local-scheme}{局部}的，可假定 $Y_{i_0}$ \Term{affine}{仿射}，因而各 $Y_i$ 与各 $X_i$ 都\Term{quasi-compact}{拟紧}且\Term{quasi-separated}{拟分离}。此时满足第 VI 讲 8.7.3 的适用条件（或者，利用 $\mathrm R^nf_*$ 作为某个\Term{presheaf}{预层}的\Term{associated-sheaf}{伴随层}来计算，可归结到 5.8）。同样可证明：
\end{proof}

\begin{SGAresult}{推论}{5.12}
断言与 5.11 相同，但现在 $F_{i_0}$ 表示 $X_{i_0}$ 上\Term{locally-finite-presentation}{局部有限表现}的交换\Term{group-scheme}{群概形}，而 $F_i,F$ 为它的\Term{inverse-image}{逆像}（\Term{scheme}{概形}\Term{base-change}{基变换}意义下的\Term{inverse-image}{逆像}）。
\end{SGAresult}
须注意，5.9（相应地，5.12）不是 5.8（相应地，5.11）的特殊情形，参见第 2 节 a）注记 2.1。

\begin{SGAresult}{推论}{5.13}
采用 5.11 的记号，设 $F$ 为 $X$ 上的\Term{abelian-sheaf}{阿贝尔群层}。则\Term{canonical-homomorphism}{典范同态}
\[
 \varinjlim_i v_i^*(\mathrm R^nf_{i*}(u_{i*}(F)))\longrightarrow\mathrm R^nf_*(F)
\]
都是\Term{isomorphism}{同构}，其中 $u_i:X\to X_i$ 和 $v_i:Y\to Y_i$ 为\Term{canonical-homomorphism}{典范同态}。
\end{SGAresult}
\SourcePage{251}
\begin{proof}
证明本质上与前面相同。
\end{proof}

\textbf{注记 5.14。} a）上述取\Term{limit}{极限}的结果，对于\Term{sheaf-of-sets}{集合层}（代替\Term{abelian-sheaf}{阿贝尔层}）的 $\mathrm H^0$ 及相应的 $f_*$ 仍成立；对于（未必交换的）\Term{sheaf-of-groups}{群层}的 $\mathrm H^1$ 及相应的 $\mathrm R^1f_*$ 也成立，其中非交换的 $\mathrm H^1$ 与 $\mathrm R^1f_*$ 按通常方式以\Term{torsor}{挠子}（或\Term{principal-homogeneous-bundle}{主齐性纤维丛}）定义（参见 Giraud 的学位论文\footnote{见原书第 235 页脚注（$*$）所引文献。}）。这可在\SGAsourcegap{原文在“dans le contexte général de”之后缺少引用}的一般框架中直接核实。对于 Giraud 研究的“\Term{band}{带}”（liens）的非交换 $\mathrm H^2$，肯定也能（而且无疑值得）给出相应的变体。

b）第 VI 讲为\Term{fibered-site}{纤维景}建立的取\Term{limit}{极限}结果，仅假定底\Term{category}{范畴}为\Term{filtered-category}{滤过范畴}，无需假定它来自一个\Term{filtered-preordered-set}{滤过预序集}。任意\Term{filtered-category}{滤过指标范畴}的情形，也是本节各断言的自然框架。我们之所以采用过于受限的框架，是为了能够正确引用 EGA IV；该书在取\Term{limit}{极限}的问题中（自第 8 节起）不幸假定\Term{index-category}{指标范畴}由\Term{filtered-preordered-set}{滤过预序集}定义。不过，以后将承认，EGA IV 中所使用的所有结果对任意\Term{filtered-category}{滤过指标范畴}都成立，因为该处给出的证明在此更一般情形下本质上无需改变。因此，当 $I$ 换成任意\Term{filtered-category}{滤过指标范畴}时，我们也将直接使用本节的结果，不再另作说明。
\SourcePage{252}
% 原 PDF 此页空白。
